\section{The quartic construction}\label{d2:sec:introduction}

\subsection*{Supported classes and generic restriction}

A class pushed forward from a Cartier divisor vanishes after restriction to the complement of that divisor. For a regular local ring, the question is whether such a class can retain infinite order in algebraic K-theory. The localization sequence and the projection formula express this question in terms of supported classes and their products; we use these results in the forms of \cite[V.3--V.6]{Weibel}. Our examples allow two different regulators to detect one such product. The first uses the analytic regulator and its comparison with the regulator of a number ring \cite{Borel,BunkeTamme}. The second uses the relative fiber comparison with cyclic homology \cite{AMMN,BGTTrace}.

The ring is defined by a quartic equation whose constant term is the residue characteristic. Its cyclotomic divisor splits into four components after a flat coefficient extension. The coefficients of these components are conjugates of one actual class in degree three. The resulting divisor calculation is the common starting point of both proofs.

\subsection{A rational generic kernel}

The geometric argument detects the class on every finite open that occurs in the localization defining the local ring.

\begin{theorem}\label{d2:thm:geometric}
Let
\[
 A=
 \left(
 \DtwoZZ_{(5)}[T,x,y]/
 \bigl(\Phi _5(1+y)+x^4+Txy^3\bigr)
 \right)_{(5,x,y)},
 \qquad
 \Phi _5(z)=z^4+z^3+z^2+z+1.
\]
Then \(A\) is a regular local ring of dimension two with residue field
\(\DtwoFF _5(T)\). There is an actual infinite-order class
\(c\in K_3(A)\) whose image in \(K_3(\DtwoFrac A)\) is zero.
Here \(K_3\) denotes Quillen algebraic K-theory.
\end{theorem}

The class is the pushforward of a cyclotomic coefficient from \(x=0\).
After the divisor splits, its analytic regulator is a linear combination
of four first Chern classes. A blowup calculation shows that the only
relation among these classes is the diagonal relation. It also controls
every geometric component of every divisor removed by an allowed
localization.

\subsection{Persistence after completion}

The completed argument requires a coefficient with prescribed nonzero
components at both odd characters.

\begin{theorem}\label{d2:thm:completed}
Let \(A\) be the ring in Theorem~\ref{d2:thm:geometric}, and let
\(\Dtwowh A\) be its \((x,y)\)-adic completion. Both rings are regular local
domains of dimension two, with maximal ideal \((x,y)\) and residue field
\(\DtwoFF _5(T)\). There is an actual class \(c\in K_3(A)\) such that its image
\(\Dtwowh c\in K_3(\Dtwowh A)\) has infinite order and
\[
 c|_{\DtwoFrac A}=0,
 \qquad
 \Dtwowh c|_{\DtwoFrac\Dtwowh A}=0
\]
integrally. Consequently both maps
\[
 K_3(A)\otimes\DtwoQQ\longrightarrow K_3(\DtwoFrac A)\otimes\DtwoQQ,
 \qquad
 K_3(\Dtwowh A)\otimes\DtwoQQ\longrightarrow
 K_3(\DtwoFrac\Dtwowh A)\otimes\DtwoQQ
\]
have nonzero kernel.
\end{theorem}

The class in Theorem~\ref{d2:thm:completed} can also be used in the geometric
proof of Theorem~\ref{d2:thm:geometric}. We establish this identification
after constructing its arithmetic coefficient. In particular, the
passage to completion does not require choosing a second supported
class.

\subsection*{Proof architecture}

The completed detector is a continuous two-form relative to a ramified
coefficient ring. Its parameter differential is retained. A tame trace
to a constant ordinary elliptic curve, followed by a bounded contraction,
takes the weighted root-divisor class to
\(4e_3(b)\) times a nonzero one-form class on a generic Cohen completion.
The exact linear identity is proved in
Proposition~\ref{d2:prop:weighted}; it holds for every anti-invariant
coefficient \(b\). Nonvanishing of the one-form rests on scalar
Frobenius pole removal in Lemma~\ref{d2:lem:scalar} and on the
rigid-to-convergent results of \cite{CaroDAddezio2025}.

Proposition~\ref{d2:prop:cyclic} identifies the relevant relative K-theory
fiber with ordinary cyclic homology after spectral completion and
inversion of \(5\). Its product calculation detects an actual relative
class. The smooth-surface finite-coefficient bound of
\cite[Theorem~8.4]{GeisserLevineCharP} then applies twice: first to the
fiber of reduction, and then to the support term in localization.
These are distinct fiber sequences.

Section~\ref{d2:sec:preliminaries} fixes the K-theory conventions and the
supported construction. Sections~\ref{d2:sec:geometry}
and~\ref{d2:sec:analytic} prove the geometric route.
Section~\ref{d2:sec:cyclic} establishes the relative cyclic comparison.
Section~\ref{d2:sec:cohen} develops the Cohen and elliptic arguments.
Section~\ref{d2:sec:quartic} computes the weighted quartic detector.
Section~\ref{d2:sec:arithmetic} constructs its actual arithmetic coefficient.
Section~\ref{d2:sec:assembly} proves Theorem~\ref{d2:thm:completed} and identifies
the coefficient used in both routes.

\section{Preliminaries}\label{d2:sec:preliminaries}

We use algebraic K-theory, transfers, localization, and finite
coefficients with the conventions of
\cite[Chapters~IV and~V]{Weibel}.

\begin{notation}\label{d2:notation:completion}
For a ring \(H\) and an ideal \(I\), put
\[
 K(H,I)=\Dtwofib\bigl(K(H)\longrightarrow K(H/I)\bigr),
 \qquad
 \DtwocK(H)=K(H)^\wedge_5[1/5],
 \qquad
 \DtwocK(H,I)=K(H,I)^\wedge_5[1/5].
\]
The completion is the derived inverse limit of the spectra
\(K(H)/5^n\), and likewise for a relative spectrum.
It is not defined by tensoring the actual K-groups with \(\DtwoZZ _5\).
Completion and inversion of \(5\) preserve fiber sequences.

On adically completed rings, differential forms are continuous forms.
Every de Rham quotient in this article uses the ordinary image of the
differential, unless a closure is explicitly introduced in the
construction of a bounded functional.
\end{notation}

For regular noetherian rings of finite dimension, projective-module
K-theory agrees with coherent-module G-theory. Finite closed immersions
have G-theory pushforwards by restriction of scalars. These pushforwards commute with flat base change and satisfy the projection formula
\[
 i_*(i^*\alpha\cdot z)=\alpha\cdot i_*z.
\]
For an effective Cartier divisor with invertible ideal \(I\), its
two-term resolution gives \(i_*[\DtwocO]=1-[I]\).
We use \cite[V.3.3, V.3.3.2, V.3.7.2, V.3.12]{Weibel}
for these statements. In each application, the immersions are finite
between affine schemes and the base change in question is flat.
Thus the required properness and Tor-independence hold.

For a regular noetherian ring \(H\), a nonzero divisor \(s\), and a
regular quotient \(H/sH\), localization and d\'evissage give
\[
 K(H/sH)\longrightarrow K(H)\longrightarrow K(H[1/s]),
\]
where the first map is support pushforward
\cite[V.4.1, V.6.1]{Weibel}.
K-groups commute with filtered colimits of rings
\cite[IV, Section~6.4]{Weibel}. The same holds with finite coefficients,
by the cofiber construction. In particular, an equality in a filtered
localization holds after some finite further localization.

\begin{setup}\label{d2:setup:ring}
Write
\[
 \begin{split}
 P(y)&=\Phi _5(1+y)=5+10y+10y^2+5y^3+y^4,\\
 F&=P(y)+x^4+Txy^3,\\
 R_0&=\DtwoZZ _{(5)}[T,x,y]/(F),\qquad
 \Dtwofm=(5,x,y),\qquad
 S=R_0\setminus\Dtwofm,\qquad A=S^{-1}R_0.
 \end{split}
\]
Fix a primitive fifth root \(\rho\), and put
\[
 E_{\Dtwoar}=\DtwoQQ(\rho),\qquad
 W_{\Dtwoar}=\DtwoZZ _{(5)}[\rho],\qquad
 \pi=\rho-1,\qquad
 G=(\DtwoZZ/5\DtwoZZ)^\times.
\]
For \(g\in G\), let \(\sigma_g(\rho)=\rho^g\).
\end{setup}

The two coefficients \(xy^3\) and \(P(y)+x^4\) of the linear polynomial
\(F\) in \(T\) have no common irreducible factor in
\(\DtwoZZ _{(5)}[x,y]\). Neither \(x\) nor \(y\) divides \(P(y)+x^4\).
Gauss's lemma therefore shows that \(R_0\) is a domain.
Solving for \(T\) gives \(\DtwoFrac A=\DtwoQQ(x,y)\).

The ambient local ring
\(\DtwoZZ _{(5)}[T,x,y]_{(5,x,y)}\) is regular of dimension three,
with residue field \(k_0=\DtwoFF _5(T)\).
The class of \(F\) in its cotangent space is the nonzero class of \(5\).
Thus \(A\) is regular local of dimension two. The identity
\begin{equation}\label{d2:eq:ramification}
 5(1+2y+2y^2+y^3)=-(x^4+y^4+Txy^3)
\end{equation}
shows that its maximal ideal is \((x,y)\).
Completion preserves the dimension and cotangent space of a noetherian
local ring, so \(\Dtwowh A\) is also regular local of dimension two, with the
same residue field. Both rings are domains.

The polynomial \(P\) is Eisenstein at \(5\).
Consequently \(W_{\Dtwoar}\) is finite free of rank four over
\(\DtwoZZ _{(5)}\). Its reduction modulo \(5\) is
\(\DtwoFF _5[\pi]/(\pi^4)\), so it is local.
Moreover,
\begin{equation}\label{d2:eq:epsilon}
 5=\pi^4\epsilon,\qquad
 \epsilon=-\frac{1}{1+2\pi+2\pi^2+\pi^3},
 \qquad
 \epsilon\bmod\pi=-1.
\end{equation}
Thus \(W_{\Dtwoar}\) is a discrete valuation ring with uniformizer \(\pi\)
and residue field \(\DtwoFF _5\). Taking the quotient by \(x\) gives
\begin{equation}\label{d2:eq:cartier-quotient}
 A/xA=W_{\Dtwoar}[T]_{(\pi)},
 \qquad y\longmapsto\pi.
\end{equation}
This is a regular discrete valuation ring. Its completion
\(\Dtwowh A/x\Dtwowh A\) is regular as well.

\begin{definition}\label{d2:defn:supported}
Let \(i\colon\DtwoSpec(A/xA)\to\DtwoSpec A\) be the closed immersion.
For an actual class \(\beta\in K_3(W_{\Dtwoar})\), define
\[
 c(\beta)=i_*\bigl(\beta|_{A/xA}\bigr)\in K_3(A).
\]
\end{definition}

Restriction to \(A[1/x]\) kills \(c(\beta)\), because the modules in its
pushforward are supported on \(x=0\). Hence its generic image is zero
integrally. The map \(A\to\Dtwowh A\) is flat.
Flat base change identifies its completed image with the same
pushforward on \(\DtwoSpec\Dtwowh A\), so that image also vanishes over
\(\DtwoFrac\Dtwowh A\). The remaining issue in both theorems is infinite order.

\section{The divisor lattice and its finite opens}\label{d2:sec:geometry}

In this section, we prove the Picard and Betti persistence used in
Theorem~\ref{d2:thm:geometric}.

\begin{proposition}\label{d2:prop:geometry}
In Setup~\ref{d2:setup:ring}, put
\[
 B_{\Dtwoar}=A\otimes_{\DtwoZZ _{(5)}}W_{\Dtwoar},
 \qquad
 \DtwocT=\DtwoSpec B_{\Dtwoar}[1/\pi,1/y],
\]
and
\[
 \DtwocQ=\DtwoSpec E_{\Dtwoar}[T,x,y,y^{-1}]/(F),
 \qquad
 D_g=V(x,y-(\rho^g-1))\subset\DtwocQ.
\]
For \(s\in S\), write \(\DtwocQ_s=D(s)\subset\DtwocQ\).
Then \(\DtwocQ\) is a smooth geometrically integral surface,
\(\DtwocT=\DtwoSpec S^{-1}\DtwocO(\DtwocQ)\), and
\[
 \DtwoPic(\DtwocQ)=\DtwoZZ ^4/\DtwoZZ(1,1,1,1),
\]
where the standard basis maps to \(\DtwocO(D_g)\).
This calculation persists under every extension of \(E_{\Dtwoar}\), and
\(\DtwoPic(\DtwocQ)\to\DtwoPic(\DtwocT)\) is injective.

For every embedding \(E_{\Dtwoar}\to\DtwoCC\), the classes
\[
 e_g=c_1^{\Dtwotopol}(\DtwocO(D_g))
       \in H^2(\DtwocQ(\DtwoCC),\DtwoRR)
\]
have exactly the real diagonal relation. For every \(s\in S\),
restriction induces an injection
\[
 H^2(\DtwocQ(\DtwoCC),\DtwoRR)\longrightarrow H^2(\DtwocQ_s(\DtwoCC),\DtwoRR).
\]
\end{proposition}

\begin{proof}
The ring \(B_{\Dtwoar}\) is finite free over \(A\).
Its closed fiber is \(k_0[\pi]/(\pi^4)\), so it has the single maximal
ideal \(\Dtwofn=(\pi,x,y)\).
Tensoring \(A\hookrightarrow\DtwoQQ(x,y)\) with the free
\(\DtwoZZ _{(5)}\)-module \(W_{\Dtwoar}\) embeds \(B_{\Dtwoar}\) in
\(E_{\Dtwoar}(x,y)\). Thus \(B_{\Dtwoar}\) is a domain.
With \(V_{\Dtwoar}=W_{\Dtwoar}[T]_{(\pi)}\), comparison of localizations gives
\[
 B_{\Dtwoar}=
 \bigl(V_{\Dtwoar}[x,y]/(F)\bigr)_{(\pi,x,y)}.
\]
Here every polynomial whose reduction is nonzero modulo \(\pi\)
becomes a unit. Inverting \(\pi\) and \(y\) gives
\(S^{-1}E_{\Dtwoar}[T,x,y,y^{-1}]/(F)\), as asserted.

Consider the blowup
\(\DtwocY=\operatorname{Bl}_{\Dtwofn}(\DtwoSpec B_{\Dtwoar})\).
Its tangent coordinates \((U:X:Y)\) correspond to \((\pi,x,y)\).
The associated graded ring is
\[
 k_0[U,X,Y]/(X^4+Y^4+TXY^3-U^4).
\]
Indeed \(F\) has order four in the ambient regular local ring, and its
nonzero initial form generates the initial ideal of the principal
ideal \((F)\). The exceptional scheme is the plane curve
\[
 C_{\mathrm{exc}}\colon X^4+Y^4+TXY^3=U^4.
\]

A geometric singular point of this curve would satisfy \(U=0\).
It cannot have \(Y=0\), since then \(X=0\).
Putting \(s=X/Y\), the remaining equations are
\[
 s^4+Ts+1=0,\qquad 4s^3+T=0.
\]
In characteristic five these imply \(T=s^3\), \(s^4=2\), and
\(T^4=3\), contrary to the transcendence of \(T\).
Thus \(C_{\mathrm{exc}}\) is geometrically smooth.
Distinct geometric plane components would intersect by B\'ezout and
give a singular point; repeated components would also be singular.
Therefore the curve is geometrically integral.

The strict-transform equations retain the coefficient
\(\epsilon=5/\pi^4\) from \eqref{d2:eq:epsilon}:
\begin{align*}
 h_\pi&=X^4+Y^4+TXY^3+
       \epsilon(1+2\pi Y+2\pi^2Y^2+\pi^3Y^3),\\
 h_x&=1+Y^4+TY^3+
       \epsilon U^4(1+2xY+2x^2Y^2+x^3Y^3),\\
 h_y&=X^4+1+TX+
       \epsilon U^4(1+2y+2y^2+y^3).
\end{align*}
The substitutions are respectively
\((x,y)=(\pi X,\pi Y)\),
\((\pi,y)=(xU,xY)\), and
\((\pi,x)=(yU,yX)\).
The latter two ambient charts retain the relations
\(\pi=xU\) and \(\pi=yU\).

These ambient charts are regular.
For example, on the \(x\)-chart one can eliminate \(x\) or \(U\) when
the other is a unit. At a point where both vanish,
\(\pi-xU\) has the nonzero linear term \(\pi\) in a regular local
polynomial ring over \(V_{\Dtwoar}\).
The argument on the \(y\)-chart is identical with \(x\) replaced by
\(y\); the \(\pi\)-chart is a localization of \(V_{\Dtwoar}[X,Y]\).

Let \(t\) denote the exceptional parameter on any one of these charts.
It is a nonzero divisor, the ambient quotient modulo \(t\) is a
polynomial domain over \(k_0\), and the corresponding \(h\) has
nonzero reduction. If \(ta=hb\), reduction modulo \(t\) forces
\(b=tb'\), and cancellation gives \(a=hb'\).
Thus the quotient by \(h\) has no \(t\)-torsion.
The displayed equations describe the strict transform without an
additional exceptional component.

At a point of \(C_{\mathrm{exc}}\), its smoothness gives a nonzero
differential of the reduced equation in an exceptional-plane
direction. Therefore the strict transform is regular there.
Away from the exceptional divisor, where \(\pi\) is invertible, the
characteristic-zero hypersurface is smooth: if
\(F_T=xy^3=0\), then either \(x=0\), where \(P'(y)\ne0\) at each
root of \(P\), or \(y=0\), where \(x^4=-5\) and
\(F_x=4x^3\ne0\).
At a nonclosed point over \(\pi=0\), the binary form
\(x^4+y^4+Txy^3\) is geometrically square-free by the repeated-root
calculation above. Its affine zero locus is smooth away from
\((x,y)=(0,0)\), so its \(x,y\) Jacobian proves regularity there.
Consequently \(\DtwocY\) is regular everywhere and is an isomorphism off
\(\Dtwofn\). The hypersurface \(B_{\Dtwoar}\) is Cohen--Macaulay and regular
in codimension one, hence normal.

The exceptional ideal satisfies
\[
 \Dtwofn\DtwocO_{\DtwocY}=\DtwocO_{\DtwocY}(-C_{\mathrm{exc}}),
 \qquad
 \DtwocO_{\DtwocY}(C_{\mathrm{exc}})|_{C_{\mathrm{exc}}}
       =\DtwocO_{C_{\mathrm{exc}}}(-1).
\]
These identities follow from the exceptional parameters on the charts,
or from \cite[Lemmas~31.33.2 and~31.33.4, Tags~0804 and~02OS]{StacksBlowup}:
the charts are \(R[I/a]\), the extended ideal is generated by \(a\),
and the degree-minus-one tautological line is the exceptional-divisor
line.

Write \(a_g=\rho^g-1\) and \(b_g=a_g/\pi\).
Then \(b_g\) is a unit with reduction \(g\).
The strict transform \(\Dtwowt D_g\) of the prime
\((x,y-a_g)\) meets the exceptional curve only at
\[
 P_g=[1:0:g].
\]
Its parameter is \(\pi\), giving this tangent direction.
On the \(\pi\)-chart, setting \(X=0\) gives
\[
 h_\pi(0,Y)=\prod_{g=1}^4(Y-b_g),
\]
because \(P(\pi Y)=\pi^4\prod_g(Y-b_g)\).
The \(Y\)-derivative is a unit at
\((\pi,X,Y-b_g)\).
The maximal ideal there is generated by \(\pi,X\), while
\(\Dtwowt D_g\) is given by \(X=0\).
Its intersection with \(C_{\mathrm{exc}}\), given by \(\pi=0\),
has multiplicity one.

We next compute the lattice on
\(C_{\mathrm{exc}}^*=C_{\mathrm{exc}}\cap\{UY\ne0\}\).
Extend the constants to \(k=\overline{\DtwoFF}_5\), and consider
\[
 M=\DtwoSpec k[X,Z,Y,Y^{-1}]/(XZ-(1-Y^4)),
 \qquad
 T=\frac{Z-X^3}{Y^3}.
\]
This surface is smooth, since the derivative \(4Y^3\) is a unit.
The open \(X\ne0\) has coordinate ring
\(k[X^{\pm1},Y^{\pm1}]\), which is factorial.
The omitted prime divisors are
\(L_g=V(X,Y-g)\), each with multiplicity one in \(\Dtwodivisor(X)\).
They generate \(\DtwoCl(M)=\DtwoPic(M)\) by divisor localization.
A relation among them is the divisor of a function that becomes a unit
on the Laurent open. Such a unit is \(cX^rY^s\).
Since \(Y\) is already invertible on \(M\), the only relation is
\[
 \DtwoPic(M)=\DtwoZZ^4/\DtwoZZ(1,1,1,1).
\]

For every \(t\in k\), the fiber \(T=t\) is
\[
 X^4+Y^4+tXY^3=1,\qquad Y\ne0.
\]
With \(s=X/Y\) and \(w=1/Y\), its coordinate ring is
\[
 k[s,w,w^{-1}]/\bigl(w^4-(s^4+ts+1)\bigr).
\]
The polynomial \(s^4+ts+1\) is not a square in \(k(s)\).
A rational square that is polynomial is a polynomial square by unique
factorization. Comparing it with \((s^2+bs+c)^2\), after normalizing
the leading coefficient, gives \(b=c=0\) from the cubic and quadratic
coefficients, contradicting the constant term.
Since \(k\) contains the fourth roots of unity and has characteristic
different from two, adjoining a fourth root of this nonsquare has
degree four. The extension is a splitting field with Galois group
inside \(\mu_4\); a subgroup of order at most two would fix the square
of the root and make the polynomial a square.
Thus every fiber is integral and reduced.
Its single vertical prime divisor is \(\Dtwodivisor(T-t)\), with
multiplicity one.

Localization to \(k(T)\) removes exactly these vertical prime
divisors. Their classes are zero, so
\[
 \DtwoPic(M)\xrightarrow{\ \sim\ }\DtwoPic(M_{k(T)}).
\]
The generic fiber is \(C_{\mathrm{exc},k(T)}^*\), and \(L_g\)
restricts to \(P_g\) with multiplicity one.
Hence a principal combination \(\sum_g n_gP_g\) on
\(C_{\mathrm{exc}}^*\) must be diagonal: extend to \(k(T)\) and use
the lattice just computed.

The substitution \(v=x^3+Ty^3\) identifies
\[
 \DtwocQ=\DtwoSpec E_{\Dtwoar}[x,v,y,y^{-1}]/(xv+P(y)).
\]
This surface is smooth because \(P\) has four distinct roots, all in
\(E_{\Dtwoar}\). Its open \(x\ne0\) is the Laurent factorial surface
with coordinates \(x,y\).
The removed prime divisors are exactly \(D_g\), each occurring once
in \(\Dtwodivisor(x)\).
Every relation is obtained from a Laurent unit \(cx^ry^s\).
It follows that
\(\DtwoPic(\DtwocQ)=\DtwoZZ^4/\DtwoZZ(1,1,1,1)\).
This also proves geometric integrality.
The same argument works over every extension of \(E_{\Dtwoar}\).

Suppose that \(\sum_g n_gD_g\) is principal on \(\DtwocT\), with
rational function \(f\).
On the regular surface \(\DtwocY\), its divisor differs from
\(\sum_g n_g\Dtwowt D_g\) by
\[
 mC_{\mathrm{exc}}+\sum_jm_jV_j,
\]
where \(V_j\) are strict transforms of primes in \(\pi=0\) or \(y=0\).
Indeed \(\DtwocY[1/\pi,1/y]=\DtwocT\).
The strict transforms over \(\pi=0\) meet the exceptional curve only
in \(U=0\), and those over \(y=0\) meet it only in \(Y=0\), as the
chart equations show.
They therefore miss \(C_{\mathrm{exc}}^*\).
The line \(\DtwocO_{C_{\mathrm{exc}}}(-1)\) is trivial where \(U\ne0\).
Restriction of the principal-divisor identity gives
\[
 \sum_g n_g[P_g]=0\quad\text{in }\DtwoPic(C_{\mathrm{exc}}^*).
\]
The coefficients \(n_g\) are equal.
The original combination is already a multiple of \(\Dtwodivisor(x)\)
on \(\DtwocQ\), which proves Picard injectivity.

Fix an embedding into \(\DtwoCC\).
The divisors \(D_g(\DtwoCC)\) are disjoint affine lines, and their
complement is \((\DtwoCC^\times)^2\).
The Thom isomorphism gives the exact segment
\[
 H^1((\DtwoCC^\times)^2,\DtwoRR)
 \longrightarrow \bigoplus_g\DtwoRR
 \longrightarrow H^2(\DtwocQ(\DtwoCC),\DtwoRR).
\]
The second map sends the standard vectors to \(e_g\).
The angular classes of \(x,y\) form a basis of the first group.
Around each \(D_g\), \(x\) winds once because its zero is simple,
whereas \(y\) winds zero times because it is invertible.
The image of the first map is precisely the diagonal line.
This proves the asserted relation among the \(e_g\).

Let \(Z\) be an irreducible \(E_{\Dtwoar}\)-divisor in the zero locus of
some \(s\in S\).
Its restriction to \(\DtwocT\) is empty.
Picard injectivity therefore gives \([Z]=0\) on \(\DtwocQ\).
Over an algebraic closure, its geometric components
\(Z_1,\ldots,Z_d\) form one Galois orbit and each has multiplicity
one, since the characteristic is zero.
The geometric Picard group is the torsion-free lattice computed
above. Galois acts trivially on it because its generators \(D_g\)
are already defined over \(E_{\Dtwoar}\).
Thus all \([Z_j]\) equal a class \(h\), and \(dh=0\).
It follows that every \([Z_j]\) vanishes.
Geometric irreducibility persists under extension of algebraically
closed fields, so every irreducible complex component of the boundary
has zero Picard class and zero topological first Chern class.

Let \(H\) be the reduced complex divisor cut out by \(s\).
If it is empty, restriction is the identity.
Otherwise remove the finite set \(\Sigma\) of singular points and
component intersections.
Local cohomology at a point of a smooth real four-manifold is
concentrated in degree four, by the cohomology of a ball relative to
its punctured ball. Therefore
\[
 H^2_\Sigma(\DtwocQ(\DtwoCC),\DtwoRR)
 =H^3_\Sigma(\DtwocQ(\DtwoCC),\DtwoRR)=0.
\]
Put \(X=\DtwocQ(\DtwoCC)\), \(X^\circ=X\setminus\Sigma\), and
\(D=H\setminus\Sigma\). The support exact sequences give
\[
 H^2_H(X,\DtwoRR)\xrightarrow{\sim}H^2_D(X^\circ,\DtwoRR),
 \qquad
 H^2(X,\DtwoRR)\xrightarrow{\sim}H^2(X^\circ,\DtwoRR).
\]
Each irreducible component \(H_i\) of \(H\) has connected smooth
part \(D_i=H_i\setminus\Sigma\): its normalization is a connected
Riemann surface, and removing finitely many points preserves
connectedness. The \(D_i\) are disjoint closed smooth curves in
\(X^\circ\).

A tubular neighborhood is a disk neighborhood in their normal
complex line bundles; a varying radius accommodates noncompact
curves. Excision reduces cohomology with supports to the relative
cohomology of this neighborhood and its complement of the zero
section. Over a contractible trivializing open, the normal fiber
pair \((B^2,B^2\setminus\{0\})\) has relative cohomology
\(\DtwoRR\) in degree two and zero otherwise.
Changes of trivialization preserve its complex orientation and
therefore its generator. Relative Mayer--Vietoris on a good
trivializing cover gives
\[
 H^2_D(X^\circ,\DtwoRR)=H^0(D,\DtwoRR)
       =\bigoplus_i\DtwoRR.
\]
This uses ordinary cohomology of \(D\), without a compactness
assumption.

The canonical section of \(\DtwocO_X(H_i)\) is transverse to zero
along \(D_i\) on \(X^\circ\), with complex-linear normal derivative.
Pulling back the oriented fiber Thom class gives the corresponding
support generator. After forgetting supports, homotopy of this
section to the zero section identifies its class with the Euler
class of the underlying oriented real plane bundle. This equals
\(c_1(\DtwocO_X(H_i))|_{X^\circ}\): both classes are given by the
winding-number cocycle of the circle-valued transition functions,
with the complex orientation fixing the positive generator.
The isomorphism on ambient \(H^2\) above identifies its preimage
with the global class \(c_1(\DtwocO_X(H_i))\) on \(X\).
Every such class is zero by the Picard-group calculation.
The restriction sequence
\[
 H^2_H(\DtwocQ(\DtwoCC),\DtwoRR)\longrightarrow H^2(\DtwocQ(\DtwoCC),\DtwoRR)
 \longrightarrow H^2(\DtwocQ_s(\DtwoCC),\DtwoRR)
\]
now proves injectivity for every \(s\).
\end{proof}

\section{The multiplicative analytic regulator}\label{d2:sec:analytic}

The goal of this section is to prove Theorem~\ref{d2:thm:geometric} from
Proposition~\ref{d2:prop:geometry}.

\begin{proof}[Proof of Theorem~\ref{d2:thm:geometric}]
Regularity and the supported construction were established in
Section~\ref{d2:sec:preliminaries}.
Let \(\DtwocO_{E_{\Dtwoar}}\) be the integer ring of \(E_{\Dtwoar}\).
It is contained in \(W_{\Dtwoar}\): an element integral over \(\DtwoZZ\)
is integral over \(W_{\Dtwoar}\), which is integrally closed in
\(E_{\Dtwoar}\).

The field \(E_{\Dtwoar}\) has two complex-conjugate pairs of embeddings
and no real embedding. By \cite[Proposition~12.2]{Borel},
\[
 \dim_\DtwoRR\bigl(K_3(\DtwocO_{E_{\Dtwoar}})\otimes\DtwoRR\bigr)=2.
\]
The analytic regulator is injective on this real vector space by
\cite[Section~6, especially~(6.4)]{BunkeTamme}.
The comparison there with the Borel regulator is up to a nonzero
factor \(2\), which does not affect injectivity.
At a point, its weight-two target is \(\DtwoCC/\DtwoRR(2)\), where
\(\DtwoRR(p)=(2\pi i)^p\DtwoRR\).
Conjugate embeddings give conjugate values.
The interval description has point representatives
\(i^{p+1}dt\); in weight two they are purely imaginary.

Choose \(\gamma_0\in K_3(\DtwocO_{E_{\Dtwoar}})\) with nonzero regulator,
and set
\[
 \gamma=\gamma_0-\sigma_{-1}\gamma_0,\qquad
 \beta=\gamma|_{W_{\Dtwoar}}.
\]
Since \(\DtwoRR(2)\) is the real subspace of \(\DtwoCC\), conjugation negates
the unique imaginary representative of a point regulator.
Thus the regulator of \(\gamma\) is twice that of \(\gamma_0\)
and remains nonzero after restriction to \(E_{\Dtwoar}\).
The identity \(\sigma_{-1}\beta=-\beta\) holds in actual K-theory.
Put \(c=c(\beta)\).
Its generic restriction is zero by Definition~\ref{d2:defn:supported}.

On \(\DtwocQ\), define
\[
 J_g=\DtwocO_{\DtwocQ}(-D_g),\qquad
 \eta_g=1-[J_g],\qquad
 \kappa=\sum_{g\in G}(\sigma_g\beta)|_{E_{\Dtwoar}}\cdot\eta_g.
\]
The resolution of the Cartier divisor gives
\(\eta_g=[\DtwocO_{D_g}]\) in \(K_0(\DtwocQ)\).
We will detect \(\kappa\) on every \(\DtwocQ_s\).

The regulator is first defined at a finite arithmetic stage.
Choose \(N\) divisible by \(5\) and by the denominators needed for
\(s\). The roots \(a_g=\rho^g-1\) and their pairwise differences are
units over \(\DtwocO_{E_{\Dtwoar}}[1/N]\).
Indeed their norms have absolute value \(5\):
for \(a_g\) this is \(\Phi_5(1)=5\), and a difference is
\(\rho^h(\rho^{g-h}-1)\).
The scheme
\[
 X_N=\DtwoSpec
 \DtwocO_{E_{\Dtwoar}}[1/N][x,v,y,y^{-1}]/(xv+P(y))
\]
is smooth over \(\DtwocO_{E_{\Dtwoar}}[1/N]\).
At a residue point where either \(x\) or \(v\) is nonzero, one of the
corresponding partial derivatives is nonzero.
Where both vanish, \(P'(y)\ne0\), since the four roots remain distinct.
Thus \(X_N\) and \((X_N)_s\) are regular separated schemes of finite
type over \(\DtwoZZ\).
The root divisors extend as Cartier divisors: near \(y=a_g\), the
other factors of \(P\) are units and the equation expresses
\(y-a_g\) as a multiple of \(x\).
The coefficient classes \(\sigma_g\gamma\) extend there as well.

For such schemes, \cite[Theorem~2.31 and Remark~2.27]{BunkeTamme}
gives a natural regulator of ring spectra with homotopy maps
\[
 K_n(X)\longrightarrow
 \prod_{p\ge0}H_{\mathcal D,\Dtwoan}^{\,2p-n}
                  (X(\DtwoCC),\DtwoRR(p)).
\]
The target satisfies the arithmetic conjugation invariance.
After constructing this map, project to the component selected by
a fixed embedding \(\tau\colon E_{\Dtwoar}\to\DtwoCC\).
Increasing \(N\) is compatible with this construction.
Filtered-colimit continuity
\cite[IV, Section~6.4]{Weibel} therefore defines the regulator on
\(K_3(\DtwocQ_s)\); the complex component is always
\(U_s=\DtwocQ_s(\DtwoCC)\).

We use its multiplicative structure explicitly.
By \cite[Remark~2.6 and Definition~2.7]{BunkeTamme}, analytic
Deligne cohomology is the shifted cone of
\[
 (2\pi i)^pA_\DtwoRR^\bullet(U)
       \oplus F^pA_\DtwoCC^\bullet(U)
 \longrightarrow A_\DtwoCC^\bullet(U),
 \qquad (a,b)\longmapsto a-b.
\]
Here \(A^\bullet\) denotes smooth forms, and \(F^p\) means at least
\(p\) holomorphic differentials. There is no logarithmic boundary
condition. Its interval model consists of forms on
\([0,1]\times U\) whose endpoint values are respectively real-twisted
and in \(F^p\).
Multiplication is wedge product.
The map to the cone is
\[
 \theta\longmapsto
 \left(\theta|_0,\theta|_1,-\int_{[0,1]}\theta\right);
\]
see \cite[Proposition~2.8 and Corollary~2.14]{BunkeTamme}.

For a complex surface, \(F^3A_\DtwoCC^\bullet(U)=0\), because at most two
holomorphic coordinate differentials can occur.
The cone sequence consequently identifies
\begin{equation}\label{d2:eq:analytic-quotient}
 H_{\mathcal D,\Dtwoan}^{3}(U,\DtwoRR(3))
   =H^2(U,\DtwoCC)/H^2(U,\DtwoRR(3)).
\end{equation}
The next possible term is the kernel of
\(H^3(U,\DtwoRR(3))\to H^3(U,\DtwoCC)\), which is zero by extension of
scalars. This argument does not require properness.

Let \(r_g\in i\DtwoRR\) be the imaginary representative of the point
regulator of \(\sigma_g\beta\) at \(\tau\).
At least one \(r_g\) is nonzero, since the embeddings
\(\tau\sigma_g\) exhaust all embeddings of \(E_{\Dtwoar}\).
Antisymmetry and commutativity of \(G\) give \(r_{-g}=-r_g\).
The vector \((r_g)\) is therefore nonzero and not diagonal.

The interval representative of this point class is \(-r_gdt\).
Choose a Hermitian metric and Chern connection on \(\DtwocO(D_g)\),
and let \(\omega_g\) be the resulting real closed representative of
\(e_g\). Its weight-one regulator is the interval-constant form
\(2\pi i\omega_g\).
This follows from the unitary, type-\((1,1)\) curvature and the
degree-two term of the unnormalized character
\(\DtwoTr\exp(-\mathrm{curvature})\) in
\cite[Definitions~2.16--2.20]{BunkeTamme}.
The weight-one class of \(\eta_g\) is the same, by additivity of
the first Chern class.
Wedge multiplication followed by minus integration gives
\[
 (-r_gdt)\wedge(2\pi i\omega_g)
 \longmapsto 2\pi i r_g\omega_g.
\]
Thus the weight-three regulator of \(\kappa|_{\DtwocQ_s}\) is
\begin{equation}\label{d2:eq:analytic-detector}
 \left[2\pi i\sum_g r_g e_g|_{U_s}\right]
 \in H^2(U_s,\DtwoCC)/H^2(U_s,\DtwoRR(3)).
\end{equation}
No other point weight contributes: point Deligne cohomology is
nonzero only in degrees zero and one, and for \(K_3\) the equation
\(2p-3=1\) forces \(p=2\).

Write \(r_g=it_g\) with \(t_g\in\DtwoRR\).
Proposition~\ref{d2:prop:geometry} shows that
\(\sum_g t_ge_g\ne0\), since \((t_g)\) is not diagonal.
It stays nonzero on every \(U_s\).
The representative in \eqref{d2:eq:analytic-detector} is the nonzero
real class
\[
 -2\pi\sum_g t_ge_g|_{U_s}.
\]
By contrast, \(H^2(U_s,\DtwoRR(3))=iH^2(U_s,\DtwoRR)\) is purely imaginary.
The real and purely imaginary subspaces intersect in zero.
Hence \eqref{d2:eq:analytic-detector} is nonzero.
Its target is a real vector space, so
\(\kappa|_{\DtwocQ_s}\) has infinite order for every \(s\in S\).

The map \(A\to B_{\Dtwoar}[1/\pi,1/y]\) is flat.
Over it, the divisor \(x=0\) splits into the disjoint components
\(D_g\cap\DtwocT\).
On the \(g\)-th component, the coefficient map sends
\(\rho=1+y\) to \(\rho^g\).
Flat base change and the projection formula give the equality of
actual classes
\begin{equation}\label{d2:eq:analytic-splitting}
 c|_{\DtwocT}
 =\sum_g(\sigma_g\beta)|_{E_{\Dtwoar}}\,[\DtwocO_{D_g\cap\DtwocT}]
 =\kappa|_{\DtwocT}.
\end{equation}
The ring of \(\DtwocT\) is regular, being a localization of the smooth
surface \(\DtwocQ\). No identification of K-theory and G-theory on the
possibly singular ring \(B_{\Dtwoar}\) is involved.

Finally,
\[
 K_3(\DtwocT)=\varinjlim_{s\in S}K_3(\DtwocQ_s).
\]
If a nonzero integer killed \(c\), it would kill \(\kappa\) in this
colimit by \eqref{d2:eq:analytic-splitting}.
The resulting equality would hold at a finite stage, contrary to
\eqref{d2:eq:analytic-detector}. Thus \(c\) has infinite order.
\end{proof}

\section{Relative cyclic detection}\label{d2:sec:cyclic}

In this section, we prove the comparison and product formula that will
detect the completed supported class.

\begin{proposition}\label{d2:prop:cyclic}
Let \(k=\overline{\DtwoFF}_5\), \(W_0=W(k)\), and let \(R\) be the
integer ring of a finite totally ramified extension of \(W_0[1/5]\).
Write \(\pi_R\) for a uniformizer and \(E_R=\DtwoFrac R\).
Let \(B_0\) be a localization of a smooth \(R\)-algebra whose relative
cotangent module is finite projective of constant rank two.
Assume that \(B=\varprojlim_nB_0/5^nB_0\ne0\).
There is a natural equivalence
\begin{equation}\label{d2:eq:relative-cyclic}
 \DtwocK(B,\pi_RB)\simeq\Sigma\DtwoHC^\Dtwocts(B/R;\DtwoQQ_5)
\end{equation}
compatible with the \(K(B)\)-module action through the negative cyclic
Chern character.

The degree-two cyclic group has the canonical summand
\[
 \Omega^2_{B/R}[1/5]\big/d\Omega^1_{B/R}[1/5].
\]
For an actual \(\beta\in K_3(R,\pi_RR)\), let \(b\in E_R\) be its
relative point coordinate, so its point character is \(t^{-1}b\)
in the mixed-HKR convention with \(|t|=-2\).
There is a scalar \(\lambda_R\in E_R^\times\), independent of
\(B\), \(J\), and \(\beta\), such that for every invertible
\(B\)-module \(J\), the top-form component of
\(\beta_B(1-[J])\) under \eqref{d2:eq:relative-cyclic} is
\[
 -\lambda_R b\,c_1^\DtwodR(J).
\]
Nonvanishing of \(-b\,c_1^\DtwodR(J)\), or of a finite sum of
these forms, proves nonvanishing of the corresponding actual
relative class after rationalized spectral completion: the whole
sum is multiplied by the same nonzero \(\lambda_R\).
No injectivity from actual K-theory to its completion is asserted.
\end{proposition}

\begin{proof}
The fiber theorem \cite[Theorem~A]{AMMN} applies to a commutative
ring henselian along \((5)\). Its coefficients \(\DtwoQQ_5\) mean
spectral completion followed by inversion of \(5\).
It gives a cartesian square with rows
\[
 K(T;\DtwoQQ_5)\longrightarrow K(T/5;\DtwoQQ_5),
 \qquad
 \DtwoHC^-(T;\DtwoQQ_5)\longrightarrow\DtwoHP(T;\DtwoQQ_5),
\]
and hence identifies the relative K-theory fiber with
\(\Sigma\DtwoHC(T;\DtwoQQ_5)\).
The continuity result
\cite[Proposition~4.2]{AMMN} applies to bounded \(5\)-power
torsion and compares completed Hochschild and cyclic objects with
their infinitesimal inverse limits.
The paragraph preceding \cite[Proposition~4.10]{AMMN}
also treats relative Hochschild and cyclic theories.
We extend that relative argument to \(B_0\) by bounds in each finite
degree range; no topological finite-presentation assertion about
\(B\) is needed.

Both \(B_0\) and \(B\) are \(5\)-torsion-free, by flatness over the
discrete valuation ring and flatness of noetherian adic completion.
Their continuous relative forms are the completions of
\(\Omega^j_{B_0/R}\). Completion is exact on these finite projective
modules, and they vanish for \(j>2\).

For every \(n\), the ring \(W_0/5^n\) is the filtered union of the
finite \'etale Witt extensions of \(\DtwoZZ/5^n\).
Its relative cotangent complex over \(\DtwoZZ/5^n\) is zero.
Flat base change therefore gives
\({L_{W_0/\DtwoZZ}}^\wedge_5=0\).
Choose an Eisenstein presentation \(R=W_0[\pi_R]\), with polynomial
\(f\). Then
\[
 {L_{R/\DtwoZZ}}^\wedge_5
   \simeq [R\xrightarrow{f'(\pi_R)}R]
\]
in cohomological degrees \(-1,0\).
Separability in characteristic zero gives \(f'(\pi_R)\ne0\).
Choose \(a\ge0\) with \(5^a\in f'(\pi_R)R\).
Multiplication by \(5^a\) on this complex is null-homotopic.

Completed cotangent transitivity gives a triangle
\[
 M\longrightarrow \Dtwowh L\longrightarrow P,
 \qquad
 M=[B\xrightarrow{f'(\pi_R)}B],\quad
 \Dtwowh L={L_{B_0/\DtwoZZ}}^\wedge_5,\quad
 P=(\Omega^1_{B_0/R})^\wedge_5.
\]
The module \(P\) is projective in degree zero and \(M\) lies in
degrees \(-1,0\).
The connecting map \(P\to M[1]\) is zero, so the triangle splits
as a triangle of \(B\)-complexes.
The kernel terms in the \(i\)-th derived exterior power of the
projection are sums of
\[
 \bigl(L\!\Lambda^rM\bigr)\otimes\Lambda^{i-r}P,
 \qquad 1\le r\le i.
\]
The \(r\)-th exterior-power functor takes the null-homotopy of
\(5^a\Dtwoid_M\) to one of \(5^{ar}\).
Thus each displayed term is killed as an object by \(5^{ar}\).

The Hochschild--Kostant--Rosenberg filtration has graded pieces
\((L\!\Lambda^iL)[i]\), and its tails become increasingly connective.
Only finitely many pieces affect any fixed homotopy range.
This remains true after derived completion.
For a \(c\)-connective spectrum, the only possible additional bottom
group is the derived inverse limit of the quotients of its
\(c\)-th homotopy group; that quotient tower is surjective.
Completion therefore preserves connectivity.
Combining the exterior-power bounds along the finitely many
extensions bounds the \(5\)-power torsion in the comparison cone in
each fixed range. Hence
\begin{equation}\label{d2:eq:hh-comparison}
 \DtwoHH(B_0;\DtwoZZ_5)[1/5]
 \xrightarrow{\ \sim\ }
 \DtwoHH(B_0/R;\DtwoZZ_5)[1/5]
\end{equation}
is an equivalence of spectra with circle action.

Ordinary cyclic homology is the circle homotopy-orbit construction.
For a connective spectrum, the omitted part of the skeletal
filtration of its bar construction becomes arbitrarily connective.
Each finite stage uses finite colimits and tensors with finite cell
spaces, which commute with spectral completion.
The finite-range argument proving \eqref{d2:eq:hh-comparison} thus
also proves the ordinary cyclic comparison.
Equivalently, the circle norm sequence makes the square comparing
the absolute and relative maps \(\DtwoHC^-\to\DtwoHP\) cartesian.
We do not need the separate vertical maps on \(\DtwoHC^-\) and \(\DtwoHP\)
to be equivalences; those constructions involve infinite limits.

The complete ring \(B\) is henselian along \((5)\).
Applying the fiber theorem to the ordinary ring \(B\), and replacing
the cyclic fiber by the relative one just computed, gives the desired
comparison with \(\DtwocK(B,5B)\).
The map \(B/5\to B/\pi_R\) has nilpotent kernel.
By \cite[Corollary~3.8]{AMMN}, a nilpotent relative extension in
a ring with nilpotent \(5\) has rationalized completed relative
K-theory zero. Exactness therefore identifies
\(\DtwocK(B,5B)\) with \(\DtwocK(B,\pi_RB)\).
This proves \eqref{d2:eq:relative-cyclic}.
Precomposing with the map from actual relative K-theory defines the
actual character.

The module structure requires more than naturality.
Put \(H=\DtwoTHH(B)\), a commutative algebra object in cyclotomic spectra.
The coefficient comparisons in
\cite[Proposition~2.8 and Corollary~2.10]{AMMN}
are tensored with \(H\); their maps are
\(\Dtwoid_H\) tensored with the coefficient maps.
They are consequently maps of \(H\)-modules.
The maps from \(H\), the linearization
\(H\otimes H\DtwoZZ\to\DtwoHH(B)\) used in
\cite[Theorem~2.12]{AMMN}, and projection to
\(\DtwoHH(B/R)\) are \(H\)-linear.

The functor \(\DtwoTC\) is lax symmetric monoidal, so these diagrams
become diagrams of \(\DtwoTC(B)\)-modules.
Circle invariants, orbits, and the norm fiber carry the corresponding
actions. The norm projection formula identifies the fiber action
with the usual \(\DtwoHC^-\)-action on \(\Sigma\DtwoHC\).
The equivalences used above may therefore be inverted in the module
category.
The relative K-to-\(\DtwoTC\) map is induced by the multiplicative
cyclotomic trace, whose \(E_\infty\) refinement is
\cite[Theorem~7.4]{BGTTrace}.
Restricting the module action along this trace proves the asserted
\(K(B)\)-linearity through the negative cyclic character.

It remains to compute the ordinary cyclic component.
In relative dimension two the normalized HKR map is integral,
because \(2\) is invertible:
\[
 a_0[a_1|\cdots|a_j]\longmapsto
 \frac1{j!}a_0\,da_1\wedge\cdots\wedge da_j
 \quad (0\le j\le2),
\]
and higher chains map to zero.
The Hochschild boundary maps to zero.
The alternating sum defining the Connes operator contributes
\(j+1\) before normalization, so the factors \(1/j!\) make it
correspond to \(d\).
For \(j=2\), the target of \(d\) is zero.
The antisymmetrized smooth HKR representative maps to the original
form, and higher relative Hochschild homology vanishes.
Thus this is a quasi-isomorphism of mixed complexes.
The completed comparison gives the finite continuous form modules
already identified.

Give a bookkeeping variable \(t\) homological degree \(-2\), and give
a \(j\)-form degree \(j\).
The mixed differential is \(td\).
The ordinary cyclic model retains the powers
\(1,t^{-1},t^{-2},\ldots\) and discards positive powers.
In degree two it has \(\Omega^2\oplus t^{-1}B\).
The differential sends \(t^{-1}b\) to \(db\).
The degree-three boundaries in the top-form term are exactly
\(d\Omega^1\).
Consequently
\begin{equation}\label{d2:eq:hc-two}
 \DtwoHC_2^\Dtwocts(B/R;\DtwoQQ_5)=
 \frac{\Omega^2_{B/R}[1/5]}{d\Omega^1_{B/R}[1/5]}
 \oplus
 t^{-1}\ker\bigl(d\colon B[1/5]\to\Omega^1_{B/R}[1/5]\bigr).
\end{equation}
In particular the image in the first summand is the ordinary
differential image.

At the coefficient point, the degree-two group is \(t^{-1}E_R\).
The class \(\beta\) has representative \(t^{-1}b\), and \(b\) is
constant for differentiation relative to \(R\).
We retain this point coordinate without further normalization.
The character comparison in
\cite[v2, Constructions~4.7--4.8 and~4.11,
Propositions~4.10 and~4.12, pp.~21--23]{AMMN}
gives a common nonzero scalar, whose image in \(E_R\) we denote
by \(\lambda_R\), for which the first Chern component of the
periodic character is \(\lambda_R c_1^\DtwodR\).
Here equation~(28) of that reference gives the special-fiber
character, equation~(24) gives the crystalline/de Rham comparison,
and equation~(25) is the fiber square.
The comparison preserves first Chern classes because their
transition functions give the same \(d\log\) cocycle.
The scalar is common to all smooth formal finite-type models
over the fixed coefficient ring; neither scalar one nor
\(\DtwoQQ_5\)-rationality is required.

This formula applies to every invertible \(B\)-module \(J\) here.
The finite presentation of \(J/\pi_RJ\), together with the finite
equations witnessing rank-one projectivity, descends to the
special fiber of a finite localization of a smooth relative
rank-two \(R\)-algebra. Complete this finite model.
An idempotent presenting the descended module lifts: starting
with a matrix lift \(e\), iterate \(e\mapsto3e^2-2e^3\).
If \(\delta=e^2-e\), the new error is
\(\delta^2(4\delta-3)\), so the errors tend to zero adically
and the matrices converge to an idempotent.
Its projective module has rank one since its reduction does
and \(\pi_R\) lies in the Jacobson radical.
The completed model maps naturally to \(B\).
After pullback, an isomorphism modulo \(\pi_R\) with \(J/\pi_RJ\)
and its inverse lift between the projective modules.
Their composites are invertible because they reduce to the
identity modulo the Jacobson radical, so the pulled-back line
is isomorphic to \(J\).
Naturality therefore carries the finite-model character formula
to \(B\).

In the rank-two mixed complex, both degree-zero negative and
periodic cyclic homology are
\[
 \ker\bigl(d\colon B[1/5]\longrightarrow
                    \Omega^1_{B/R}[1/5]\bigr)
 \ \oplus\
 t\bigl(\Omega^2_{B/R}[1/5]/d\Omega^1_{B/R}[1/5]\bigr).
\]
The comparison is the identity on these summands: total degree
zero contains functions and \(t\)-times two-forms, total degree
one supplies the boundaries \(t\,d\Omega^1\), and higher forms
vanish. Thus the negative-to-periodic map is injective in this
degree. The commutative fiber square identifies the negative
cyclic character of \(J\) with
\(1+\lambda_Rt\,c_1^\DtwodR(J)\).
The character of \(1-[J]\) is consequently
\(-\lambda_Rt\,c_1^\DtwodR(J)\), whose action on \(t^{-1}b\)
has top component \(-\lambda_Rb\,c_1^\DtwodR(J)\).
At the normalized bar level, the shuffle term gives this wedge
product, while cyclic-shuffle corrections insert the relative
scalar \(b\) in a bar slot and vanish.
This is the module action established above.
The common nonzero scalar preserves nonvanishing of every
finite sum.

A nonzero top-form component cannot be the image of a zero
rationalized completed relative class.
The same calculation is additive for finite sums.
It makes no assertion that actual K-theory injects into completion.
\end{proof}

\section{Cohen completion and elliptic contraction}\label{d2:sec:cohen}

We prove that the graph-divisor class survives Cohen completion.

\subsection{Scalar Frobenius pole removal}

\begin{lemma}\label{d2:lem:scalar}
Let \(p\) be a prime, \(k=\overline{\DtwoFF}_p\), \(W=W(k)\), and
\(K=W[1/p]\).
Let \(D\) be a smooth affine \(W\)-algebra of relative dimension one
with connected nonempty special fiber \(U=\DtwoSpec(D/pD)\).
Put
\[
 S_D=\{s\in D\mid s\bmod p\ne0\},\qquad
 \Dtwowh D=\varprojlim_nD/p^nD,\qquad
 C=\varprojlim_nS_D^{-1}D/p^nD,\qquad F_C=C[1/p].
\]
Let \(D^\dagger\) be the weak completion.
Assume that an integral Frobenius lift
\(\phi\colon D^\dagger\to D^\dagger\) restricts to the \(d\)-th Witt
Frobenius on \(W\), where \(d\ge1\), and reduces to the
\(q=p^d\) power map.
If \(\omega\in\Omega^1_{D^\dagger/W}[1/p]\) has nonzero rigid
cohomology class and
\[
 \phi^*\omega-a\omega=dh
 \qquad(a\in W,\ h\in D^\dagger[1/p]),
\]
then its image in
\[
 \Omega^1_{C/W}[1/p]\big/dF_C
\]
is nonzero.
\end{lemma}

\begin{proof}
A smooth connected curve has disjoint irreducible components, since
its local rings are domains. Therefore \(D_0=D/pD\) is an integral
normal domain. Write \(L=\DtwoFrac D_0\).
The localization \(S_D^{-1}D=D_{pD}\) is a regular local ring of
dimension one with parameter \(p\) and residue field \(L\).
Thus \(C\) is a complete discrete valuation ring with these same
data.

For every \(n\), the map
\(D/p^nD\to S_D^{-1}D/p^nD\) is injective.
If \(sa\in p^nD\) with \(s\in S_D\), reduction modulo \(p\) gives
\(a\in pD\). Write \(a=pa_1\), cancel \(p\) by flatness over \(W\),
and repeat.
Taking limits embeds \(\Dtwowh D\) in \(C\) and gives
\begin{equation}\label{d2:eq:integral-intersection}
 \Dtwowh D\cap p^nC=p^n\Dtwowh D.
\end{equation}
The continuous forms on \(C\) are the completions of the forms on
\(S_D^{-1}D\).

We first determine its continuous constants:
\begin{equation}\label{d2:eq:cohen-constants}
 \ker\bigl(d\colon F_C\to\Omega^1_{C/W}[1/p]\bigr)=K.
\end{equation}
After inverting an element of \(S_D\), smoothness supplies an
\'etale coordinate \(t\) and a ring \(D'\) \'etale over \(W[t]\).
The residue of \(t\) is a separating parameter for \(L\), and the
generic-special completion of \(D'\) is still \(C\).
Formal \'etaleness extends \(t\mapsto t+z\) to a Taylor map
\(D'\to C[[z]]\), successively modulo powers of \(z\).
Further denominators have unit constant coefficient, so this map
extends to the localization and then to \(C\), coefficient by
coefficient.
Write its coefficients as integral continuous \(W\)-linear operators
\(H_j\colon C\to C\).

Uniqueness of the \'etale lift identifies successive translations by
\(z,w\) with translation by \(z+w\).
With \(\partial=H_1\), this gives
\(j!H_j=\partial^j\) and \(\partial(t)=1\).
The continuous cotangent module is free on \(dt\).
If \(f\in C\) satisfies \(df=0\), then
\(j!H_j(f)=0\) for every \(j>0\).
The ring \(C\) is torsion-free over \(\DtwoZZ\), so \(H_j(f)=0\).

On reduction modulo \(p\), all positive Hasse operators annihilate
\(\bar f\in L\).
For each \(r\ge1\),
\[
 1,t,\ldots,t^{p^r-1}
\]
is a basis of \(L\) over \(L^{p^r}\).
Indeed finite separability of \(L/k(t)\) gives
\([L:L^p]=p\), and \(dt\ne0\) implies \(t\notin L^p\).
Iteration gives \([L:L^{p^r}]=p^r\).
A smaller purely inseparable degree for \(t\) over \(L^{p^r}\)
would force \(t\) to be a \(p\)-th power in \(L\).

Expand \(\bar f=\sum_{i=0}^{p^r-1}b_i^{p^r}t^i\).
The Taylor series of \(b_i^{p^r}\) has no terms of degrees
\(1,\ldots,p^r-1\).
Comparing the coefficients of \(z^{p^r-1},z^{p^r-2},\ldots,z\)
successively gives \(b_i=0\) for every \(i>0\).
Thus \(\bar f\in\bigcap_rL^{p^r}=k\).
For the last equality, the valuation of a nonzero element of the
intersection at every point of the smooth proper model of \(L\) is
divisible by all \(p^r\), hence zero.
Such an element is a global regular function on a proper connected
curve and belongs to \(k\); conversely \(k\) is perfect.

Choose \(w_0\in W\) lifting \(\bar f\).
Then \((f-w_0)/p\) is integral and again has zero differential.
Repeating gives \(f=\sum_{i\ge0}p^iw_i\), with \(w_i\in W\).
This converges in \(W\) and in \(C\), so \(f\in W\).
Scaling by a power of \(p\) proves \eqref{d2:eq:cohen-constants}.

The integral map \(\phi\) fixes \(p\) and extends continuously to
\(\Dtwowh D\). Its reduction preserves nonzero residues, so inversion
of \(S_D\) and completion extend it to \(C\).
We claim that
\begin{equation}\label{d2:eq:pole-removal}
 f\in C,\quad \phi(f)-af\in\Dtwowh D
 \quad\Longrightarrow\quad f\in\Dtwowh D.
\end{equation}
Write \(\phi(f)-af=r\).
A pole of \(\bar f\) on \(U\) of order \(m>0\) would give pole
order \(qm\) in \(\bar f^q\) and at most \(m\) in
\(\bar a\bar f\). These cannot cancel.
Hence \(\bar f\) has no poles on \(U\).
Normality gives \(\bar f\in D_0\).

Lift \(\bar f\) to \(g_0\in D\), and put
\[
 f_1=(f-g_0)/p,\qquad
 r_1=(r-\phi(g_0)+ag_0)/p.
\]
The numerator defining \(r_1\) belongs to
\(\Dtwowh D\cap pC=p\Dtwowh D\) by \eqref{d2:eq:integral-intersection}.
Thus \(r_1\in\Dtwowh D\) and \(\phi(f_1)-af_1=r_1\).
Iteration yields
\[
 f=\sum_{i=0}^{n-1}p^ig_i+p^nf_n,
 \qquad g_i\in D,\quad f_n\in C.
\]
The partial sums converge in \(\Dtwowh D\), and the remainder tends to
zero in \(C\). This proves \eqref{d2:eq:pole-removal}.
No \(W\)-linearity of \(\phi\) was used.

Suppose now that \(\omega=df\) in \(F_C\).
The Frobenius equation gives
\(d(\phi(f)-af-h)=0\).
By \eqref{d2:eq:cohen-constants},
\(\phi(f)-af=h+c\) for \(c\in K\).
Multiply by one sufficiently large power of \(p\) to make
\(f,h,c\) integral in \(C,\Dtwowh D,W\), respectively.
The equation is preserved because \(\phi(p)=p\).
Equation~\eqref{d2:eq:pole-removal} then puts the scaled \(f\) in
\(\Dtwowh D\). Undoing the scaling makes \(\omega\) exact in the
rational completed de Rham complex.

For the smooth finite-type curve \(U\) and its constant overconvergent
Frobenius coefficient, the map
\[
 H^1_\Dtworig(U/K)\longrightarrow H^1_\Dtwoconv(U/K)
\]
is injective by \cite[Theorem~3.2.1]{CaroDAddezio2025}.
The weakly completed de Rham complex computes the first group, and
the completed de Rham complex computes the second, with comparison
induced by \(D^\dagger\to\Dtwowh D\);
see also \cite[Theorem~2.1.2]{CaroDAddezio2025}.
Both use ordinary differential images in degree one.
The primitive just constructed contradicts this injectivity.
Only the original finite-type curve \(U\) enters this external
comparison.
\end{proof}

\subsection{A bounded elliptic contraction}

\begin{setup}\label{d2:setup:elliptic}
For the remainder of this section, let
\(k=\overline{\DtwoFF}_5\), \(W=W(k)\), and \(K=W[1/5]\).
Let
\[
 \DtwocE/W\colon v^2=z^3+z
\]
be the smooth projective elliptic curve with origin \(O\), and put
\[
 \DtwocE^\circ=\DtwocE\setminus\{O,(0,0)\},
 \qquad
 E_0=\DtwocE\otimes_Wk,
 \qquad
 A_{\mathrm{ell}}=\Gamma(\DtwocE^\circ,\DtwocO).
\]
Let \(D\) be a smooth affine \(W\)-algebra of relative dimension one
with connected nonempty special fiber. Let \(C\) be its
generic-special completion, and put
\[
 B_{\mathrm{ell}}=
 \bigl(\Dtwowh A_{\mathrm{ell}}\otimes_WC\bigr)^\wedge_5.
\]
\end{setup}

\begin{proposition}\label{d2:prop:contraction}
Choose rational dagger one-forms \(\omega,\eta\) on
\(\DtwocE^\circ\) representing nonzero generators of the proper
slope-one and slope-zero crystalline lines, respectively.
There is a bounded \(K\)-linear functional
\[
 \ell\colon\Omega^1_{\Dtwowh A_{\mathrm{ell}}/W}[1/5]\longrightarrow K
\]
such that
\[
 \ell(\omega)=1,\qquad
 \ell(\eta)=0,\qquad
 \ell(d\Dtwowh A_{\mathrm{ell}}[1/5])=0.
\]
It extends to a continuous degree-minus-one chain operator
\[
 L\colon\Omega^\bullet_{B_{\mathrm{ell}}/W}[1/5]
       \longrightarrow\Omega^\bullet_{C/W}[1/5][-1]
\]
with \(L_2(\alpha\wedge\gamma)=\ell(\alpha)\gamma\), where the
elliptic form is placed first.
The map
\[
 J\colon H^1_\DtwodR(C/W)[1/5]\longrightarrow
          H^2_\DtwodR(B_{\mathrm{ell}}/W)[1/5],
 \qquad [\gamma]\longmapsto[\omega\wedge\gamma],
\]
is split injective, with left inverse induced by \(L\).
In particular, a class represented by
\(\omega\wedge\gamma+\eta\wedge\delta\) is nonzero whenever
\([\gamma]\ne0\) in the ordinary one-form quotient of \(C\).
\end{proposition}

\begin{proof}
The cubic has distinct roots in characteristic five, so \(E_0\) is
smooth. The coefficient of \(z^4\) in \((z^3+z)^2\) is \(2\),
which is nonzero modulo five.
The Hasse-invariant criterion
\cite[Theorem~V.4.1(a), p.~148]{SilvermanAEC2009} gives ordinarity.
Proper rational crystalline \(H^1\) therefore has two lines of
slopes zero and one.
Restriction to \(E_0^\circ\) is injective in rigid \(H^1\).
The removed reduced finite set \(Z=\{O,(0,0)\}\) has codimension
one in the smooth curve \(E_0\).
The support-vanishing assertion of
\cite[Theorem~1.1, third assertion, p.~63]{Petrequin2003}
gives \(H^1_{Z,\Dtworig}(E_0/K)=0\), since \(1<2\).
The localization sequence
\cite[Proposition~1.3, p.~64]{Petrequin2003} therefore gives
\[
 0\longrightarrow H^1_\Dtworig(E_0/K)
       \longrightarrow H^1_\Dtworig(E_0^\circ/K).
\]
This natural restriction respects Frobenius.

For unramified Witt coefficients,
\cite[Theorem~7.3.1]{CaroDAddezio2025} identifies slope-one rigid and
convergent cohomology in degree one.
For smooth affine schemes with unit-root coefficients, the
slope-one convergent part meets the closure of zero trivially
\cite[Theorem~7.3.6]{CaroDAddezio2025}.
The slopes in \([0,1)\) lie in the closure of zero by
\cite[Theorem~7.2.1]{CaroDAddezio2025}.
All hypotheses apply to \(E_0^\circ\), constant coefficients, and
the unramified ring \(W\).

Put \(M_0=\Omega^1_{\Dtwowh A_{\mathrm{ell}}/W}\), \(M=M_0[1/5]\),
and let \(N\) be the closure of
\(d\Dtwowh A_{\mathrm{ell}}[1/5]\) in \(M\).
All one-forms are closed, since the relative dimension is one.
The topology and convergent-cohomology description in
\cite[Definition~7.1.1]{CaroDAddezio2025} identify the separated
quotient with \(M/N\).
The class of \(\omega\) in this quotient is nonzero, whereas
\(\eta\in N\).
We use this separated quotient only to construct \(\ell\).

The finite projective, complete, separated, torsion-free module
\(M_0\) defines a Banach norm on \(M\) with discrete values.
The quotient \(V=M/N\) is Banach for the quotient norm.
For a nonzero quotient vector, the infimum of the norms of its
representatives is positive and attained, because the norm values
are discrete and \(N\) is closed.
Completeness follows by taking a Cauchy subsequence, lifting its
successive differences with norms tending to zero, and summing
those lifts in \(M\).

Let \(V_0\) be the unit ball, and lift a basis of the \(k\)-vector
space \(V_0/5V_0\) to elements \(e_i\in V_0\).
Every \(v\in V_0\) has a unique expansion
\[
 v=\sum_i a_ie_i,\qquad a_i\in W,
\]
where the coefficients tend to zero off finite subsets.
To obtain the expansion, reduce modulo five, subtract a finite
lifted basis expansion, divide by five, and repeat.
The resulting series converges by completeness.
For uniqueness, a nonzero coefficient family has a largest absolute
value. Scaling to make that value one and reducing modulo five
contradicts independence of the finitely many largest coefficients.
Scaling also gives such expansions over \(K\) for all \(v\in V\).
The coordinate maps are consequently bounded and \(K\)-linear.

Some coordinate is nonzero on the image of \(\omega\).
Divide that coordinate by its value on \(\omega\) and compose with
\(M\to V\). This gives \(\ell\).
It annihilates \(N\), hence both \(\eta\) and exact forms.
For some \(e\ge0\), boundedness gives \(5^e\ell(M_0)\subset W\).

Set \(\ell_0=5^e\ell\).
Reduce this integral map modulo \(5^n\), tensor with \(C/5^nC\),
and pass to the inverse limit.
This defines \(M_0\Dtwoct_WC\to C\).
Tensoring instead with \(\Omega^1_{C/W}/5^n\) defines
\(M_0\Dtwoct_W\Omega^1_{C/W}\to\Omega^1_{C/W}\).
After rationalization divide by \(5^e\).
These are limits of compatible operators, with no interchange
between cohomology and inverse limits.

At each finite precision, the tensor-product formula for
differentials gives
\begin{align}
 \Omega^1_{B_{\mathrm{ell}}/W}
 &=(M_0\Dtwoct_WC)\oplus
   (\Dtwowh A_{\mathrm{ell}}\Dtwoct_W\Omega^1_{C/W}),\label{d2:eq:product-one}\\
 \Omega^2_{B_{\mathrm{ell}}/W}
 &=M_0\Dtwoct_W\Omega^1_{C/W}.\label{d2:eq:product-two}
\end{align}
Higher forms vanish.
The two factors have cotangent rank one; for \(C\), this follows
from completion of an \'etale-coordinate localization of \(D\).
The form modules are finite projective.
Flatness over \(W\) ensures that completed tensor products reduce to
the ordinary tensor products modulo \(5^n\), so the displayed limits
are valid.

Define \(L_0=0\).
Define \(L_1\) by the extended \(\ell\) on the first summand of
\eqref{d2:eq:product-one} and zero on the second.
Define \(L_2\) by the second tensor extension.
On functions, \(L_1d=0\).
For an elliptic one-form \(\alpha\) and a base function \(f\),
\[
 d(\alpha f)=-\alpha\wedge d_Cf,\qquad
 L_2d(\alpha f)=-\ell(\alpha)d_Cf=-d_CL_1(\alpha f).
\]
For an elliptic function \(a\) and a base one-form \(\gamma\),
\(d(a\gamma)=d_Aa\wedge\gamma\), so its image under \(L_2\) is
zero, as is \(d_CL_1(a\gamma)\).
Finite sums of pure tensors are dense, and all operators are
continuous. Thus \(L_2d=-d_CL_1\) on the completed complex, which is
the chain identity for the shift \([-1]\).

The form \(\omega\wedge\gamma\) is closed.
If \(\gamma=d_Cf\), then
\(\omega\wedge\gamma=-d(f\omega)\).
Hence \(J\) is defined on ordinary cohomology.
Finally \(L_2(\omega\wedge\gamma)=\gamma\).
This proves the splitting, and \(\ell(\eta)=0\) proves the last
assertion.
\end{proof}

\subsection{Graph divisors}

\begin{proposition}\label{d2:prop:graph}
In Setup~\ref{d2:setup:elliptic}, assume that
\(\Dtwowt Q\colon\DtwoSpec D\to\DtwocE^\circ\) has nonconstant special fiber
\(Q\colon U\to E_0^\circ\).
Let \(\DtwocL_Q\) be the pullback to \(\DtwoSpec B_{\mathrm{ell}}\) of
\[
 \DtwocO(\Gamma_{\Dtwowt Q}-\Gamma_{-\Dtwowt Q})
\]
on \(\DtwocE^\circ\times_W\DtwoSpec D\).
Then
\[
 c_1^\DtwodR(\DtwocL_Q)\ne0
 \quad\text{in}\quad
 \Omega^2_{B_{\mathrm{ell}}/W}[1/5]\big/
       d\Omega^1_{B_{\mathrm{ell}}/W}[1/5].
\]
Nonvanishing persists under every finite extension of \(K\), with
continuous relative forms over the extended integer ring.
\end{proposition}

\begin{proof}
Choose the proper slope-zero generator \(\eta\) and slope-one
generator \(\omega\) over \(\DtwoQQ_5\), and normalize
\(\DtwoTr(\eta\cup\omega)=1\).
The curve is defined over \(\DtwoFF_5\).
Proper smooth crystalline/de Rham comparison gives corresponding
algebraic de Rham classes on \(\DtwocE_K\).
Their restrictions to \(\DtwocE_K^\circ\) have algebraic one-form
representatives because the curve is affine.
These also define dagger and continuous representatives, so
Proposition~\ref{d2:prop:contraction} applies.

With the principal-polarization convention, the normalized
Poincar\'e bundle is
\[
 \DtwocP=\DtwocO(\Delta-\DtwocE\times O-O\times\DtwocE).
\]
Its restrictions to the origin slices are trivial, so its Chern
class has only a mixed term in the proper K\"unneth decomposition.
The diagonal acts as the identity on \(H^1\).
With our pairing convention, the mixed tensor is
\[
 \omega\otimes\eta-\eta\otimes\omega.
\]
Multiplying by \(\eta\) on the first factor and tracing returns
\(\eta\) on the second. Multiplying by \(\omega\) returns
\(\omega\). Perfectness of the pairing determines the tensor.
The trace, duality, K\"unneth, and proper crystalline compatibilities
used here are those of \cite[\S1.7(i), \S2.3(i), Theorems~2.4 and~3.2, and \S3.3(i)]{Berthelot1997Duality}.

The graph line for a map \(q\) is the pullback of the diagonal line
by \(\Dtwoid\times q\).
In the difference for \(\Dtwowt Q\) and \(-\Dtwowt Q\), the first-factor
origin line cancels, and the second-factor terms cancel because
\([-1]^*\DtwocO(O)\simeq\DtwocO(O)\).
Inversion acts by \(-1\) on proper \(H^1\).
Thus
\begin{equation}\label{d2:eq:graph-class}
 c_1^\DtwodR(\DtwocL_Q)
 =
 \left[2\bigl(\omega\wedge\Dtwowt Q^*\eta
             -\eta\wedge\Dtwowt Q^*\omega\bigr)\right].
\end{equation}
This equality first holds on the finite-type affine product over
\(K\). The difference from another Chern representative is the
differential of an algebraic one-form there.
It remains exact under weak completion, generic localization, and
continuous completion. No K\"unneth formula for the completed
product is being used.

Compactify \(U\) to its smooth proper connected curve \(\bar U\).
Properness extends \(Q\) to
\(\bar Q\colon\bar U\to E_0\).
The map \(\bar Q\) is finite and flat of rank
\(m=[k(\bar U):k(E_0)]>0\): its finite direct image is
torsion-free over the local discrete valuation rings of \(E_0\).
This rank includes the inseparable degree.
Write \(\xi_C\) for the rigid homological fundamental class of a
smooth proper curve \(C\), viewed as its top-degree trace functional.
Proper cycle pushforward gives
\(\bar Q_*[\bar U]=m[E_0]\); compatibility of cycle classes with
proper pushforward therefore gives
\(\bar Q_*\xi_{\bar U}=m\xi_{E_0}\)
\cite[Proposition~2.2, Section~2.1, and Proposition~6.4]{Petrequin2003}.
The homological pushforward is the transpose of cohomological
pullback. Consequently, for every
\(u\in H^2_{\mathrm{rig}}(E_0/K)\),
\[
 \DtwoTr_{\bar U}(\bar Q^*u)=m\DtwoTr_{E_0}(u).
\]
Taking \(u=a\cup b\) and using the naturality of cup products gives
\[
 \DtwoTr_{\bar U}(\bar Q^*a\cup\bar Q^*b)
   =m\,\DtwoTr_{E_0}(a\cup b).
\]
The pairing is perfect and \(m\) is invertible in \(K\).
Consequently \(\bar Q^*\) is injective on degree-one cohomology,
without a separability assumption.
Restriction from \(\bar U\) to \(U\) is also injective.
If the reduced finite complement \(Z=\bar U\setminus U\) is
nonempty, it has codimension one in the smooth curve \(\bar U\).
Support vanishing gives \(H^1_{Z,\Dtworig}(\bar U/K)=0\), and
the localization sequence then makes
\(H^1_\Dtworig(\bar U/K)\to H^1_\Dtworig(U/K)\) injective
\cite[Theorem~1.1, third assertion, and Proposition~1.3,
pp.~63--64]{Petrequin2003}.
For an empty complement restriction is the identity.
Thus \(Q^*\eta\ne0\) in \(H^1_\Dtworig(U/K)\), represented
by \(\Dtwowt Q^*\eta\).

The special curve and map descend to some finite field
\(\DtwoFF_{5^d}\).
The four \(\DtwoFF_5\)-points of \(E_0\) are \(O\) and the three
points with \(v=0\), at \(z=0,2,3\).
At \(z=1,4\) the cubic is a nonsquare.
The Frobenius polynomial is therefore \(X^2-2X+5\).
Its reduction has simple roots \(0,2\), so its two roots lie in
\(\DtwoZZ_5\), one a unit \(\alpha\).
We choose \(\eta\) with Frobenius eigenvalue \(\alpha\).

The weak completion \(D^\dagger\) admits an integral, Witt-semilinear
Frobenius lift by the smooth affine lifting statement in
\cite[Introduction, p.~198]{DLZ2011}.
Iterate it \(d\) times to obtain \(\phi\).
Its action on rigid cohomology is canonical and independent of the
lift.
Descent of \(Q\) gives
\(\phi^*[Q^*\eta]=\alpha^d[Q^*\eta]\).
On the affine dagger complex this means
\[
 \phi^*(\Dtwowt Q^*\eta)-\alpha^d\Dtwowt Q^*\eta=dh
 \qquad(h\in D^\dagger[1/5]).
\]
Lemma~\ref{d2:lem:scalar} applies, since \(\alpha^d\in W\).
Hence the image \([\Dtwowt Q^*\eta]\) in the ordinary Cohen one-form
quotient is nonzero.

Apply \(L\) from Proposition~\ref{d2:prop:contraction} to
\eqref{d2:eq:graph-class}.
The result is \(2[\Dtwowt Q^*\eta]\), which is nonzero.
Because \(L\) takes exact forms to exact forms, the graph class
cannot be exact.

For a finite extension \(K'/K\), its integer ring \(W'\) is finite
free over \(W\) and complete.
Scalar extension therefore commutes with the adic limits defining
the product and its finite projective form modules.
Relative forms over \(W'\) are the old forms tensored with \(W'\).
After inversion of five, the complexes are tensored with \(K'\).
Field extension is exact and faithfully flat, so it commutes with
kernels and ordinary images and preserves the nonzero class.
\end{proof}

\section{The quartic quotient and its weight}\label{d2:sec:quartic}

In this section, we construct the elliptic quotient and compute the
linear functional on the four root-divisor classes.

\begin{proposition}\label{d2:prop:weighted}
Let \(k=\overline{\DtwoFF}_5\), \(W=W(k)\), \(R=W[\rho]\), and
\[
 V=\varprojlim_n W[T]_{(5)}/5^n.
\]
There is a finite unramified extension \(C/V\), obtained as the
generic-special completion of a smooth affine \(W\)-curve \'etale
over an open of the \(T\)-line, with the following property.
Put \(C'=C\otimes_WR\) and
\[
 B=C'\langle X,Y\rangle/
 \left(X^4+\frac{\Phi_5(1+\pi Y)}{\pi^4}+TXY^3\right).
\]
For \(g\in G=\DtwoFF_5^\times\), set
\[
 a_g=\frac{\rho^g-1}{\pi},\qquad
 J_g=(X,Y-a_g),\qquad
 \gamma_g=-c_1^\DtwodR(J_g).
\]
The ideals \(J_g\) are invertible.
Let \(E=\DtwoQQ_5(\rho)\), let
\(\omega_T\colon G\to\DtwoZZ_5^\times\) be the Teichm\"uller character,
and put
\[
 e_r=\frac14\sum_{g\in G}\omega_T(g)^{-r}\sigma_g
 \qquad(r=1,3)
\]
on \(\DtwoQQ_5\)-vector-space targets.
If \(b\in E\) satisfies
\(\sigma_{-1}b=-b\) and \(e_1b,e_3b\ne0\), then
\[
 \sum_{g\in G}\sigma_g(b)\gamma_g\ne0
 \quad\text{in}\quad
 \Omega^2_{B/R}[1/5]\big/d\Omega^1_{B/R}[1/5].
\]
The relative forms retain \(dT\).
\end{proposition}

\begin{proof}
Over \(k(T)\), consider
\[
 H\colon X^4+Y^4+TXY^3=Z^4,\qquad
 H^\circ=H\cap\{Z\ne0\},\qquad
 P_g=[0:g:1].
\]
The singular-point calculation in
Proposition~\ref{d2:prop:geometry} applies with \(Z\) in place of \(U\).
It gives \(T^4=3\) at a hypothetical singularity.
Thus \(H\) is smooth and geometrically integral.

The same auxiliary surface
\[
 \DtwoSpec k[X,U,Y,Y^{-1}]/(XU-(1-Y^4)),
 \qquad T=(U-X^3)/Y^3,
\]
has root-divisor lattice \(\DtwoZZ^4/\DtwoZZ(1,1,1,1)\).
For clarity, the mechanisms remain unchanged here: inversion of \(X\)
gives the factorial Laurent ring; its units give only the diagonal
relation; each fiber becomes
\(w^4=s^4+ts+1\), whose right side is not a square by comparison
with \((s^2+bs+c)^2\); hence each vertical divisor is the reduced
principal divisor \(T-t\).
Divisor localization therefore gives this same lattice on
\(H^\circ\cap\{Y\ne0\}\).
A torsion relation on \(P_1-P_{-1}\) in the proper Jacobian would
restrict to a torsion relation in this lattice, which is impossible.
Finite field extension preserves nontorsion: norm after pullback multiplies divisor classes by the extension degree.

We construct the quotients over an \'etale parameter cover.
Set \(\Delta=256-27T^4\), and order the roots
\(b_1,\ldots,b_4\) of \(s^4+Ts+1\) over a finite \'etale cover of
\(W[T,1/\Delta]\).
Put
\[
 a=\frac{b_2-b_4}{b_2-b_1},\qquad
 x=a\frac{s-b_1}{s-b_4},\qquad
 \lambda=a\frac{b_3-b_1}{b_3-b_4}.
\]
The elements \(a,a-1,a-\lambda,\lambda,\lambda-1\) are units,
as follows from the pairwise root differences.
Adjoin \(c\) with \(c^4=a(a-1)(a-\lambda)\).
Replacing \(w=Z/Y\) by \(cw/(s-b_4)\) gives
\begin{equation}\label{d2:eq:cyclic-quartic}
 w^4=x(x-1)(x-\lambda).
\end{equation}
Indeed the product of the three transformed factors is
\(a(a-1)(a-\lambda)\prod_{j=1}^3(s-b_j)/(s-b_4)^3\),
which is \(c^4(s^4+Ts+1)/(s-b_4)^4\).

Adjoin \(r\) with \(r^2=\lambda\) and
\(\alpha_\pm\) with \(\alpha_\pm^2=r\pm1\).
All these extensions are \'etale, because the radicands and \(2,4\)
are units.
Set
\[
 z=\frac{w^2}{x},\qquad v_\pm=\frac{w(x\pm r)}x.
\]
Equation~\eqref{d2:eq:cyclic-quartic} gives
\begin{equation}\label{d2:eq:elliptic-quotients}
 z^2=x+\frac{r^2}{x}-1-r^2,\qquad
 v_\pm^2=z^3+(r\pm1)^2z.
\end{equation}
After \(z=\alpha_\pm^2\zeta\), \(v_\pm=\alpha_\pm^3v\),
the target is the constant elliptic curve \(v^2=\zeta^3+\zeta\).

Let \(\tau(x,w)=(x,-w)\) and
\[
 \iota(x,w)=\left(\frac{r^2}{x},\frac{rw}{x}\right).
\]
The maps \(h_+,h_-\) are the quotients by \(\iota\) and
\(\tau\iota\), respectively.
Their degree is two: \(x\) satisfies a quadratic over the target
functions by \eqref{d2:eq:elliptic-quotients}, \(w\) is recovered from
\(x,v_\pm\), and the displayed nontrivial involution fixes the
target functions.
They extend between the proper smooth curves.

The boundary \(Z=0\) consists of the four points
\(x=0,1,r^2,\infty\).
At \(0,\infty\), the orders of \(z,v_\pm\) are \((-2,-3)\), so
the image is \(O\) with local degree one.
At \(1,r^2\), these orders are \((2,1)\), so the image is
\((0,0)\), again with local degree one.
Everywhere else \(z\) is finite and nonzero.
The exact inverse image of \(\{O,(0,0)\}\) is therefore the boundary.
Both maps restrict to finite maps of all of \(H^\circ\) to
\(E_0^\circ\).

The equations are defined over \(W\).
Shrinking the parameter open by elements nonzero at its generic
special point makes these finite flat quotient maps between smooth
relative curves.
Choose a connected component of the special fiber of the finite
\'etale parameter cover.
It can be isolated in a smooth affine \(W\)-model by inverting a
lift of its characteristic idempotent.
Write \(D\) for the resulting smooth \(W\)-algebra.
Its generic-special completion \(C\) is the corresponding finite
unramified extension of \(V\); subsequent open restrictions are
units in \(C\).

The involutions \(\tau,\iota,\tau\iota\) form a Klein four group.
Its total quotient is rational with coordinate \(z=w^2/x\).
Recovering \(x\) and then \(w\) gives degree at most four, and the
four automorphisms fix \(z\), so \(k(T)(z)\) is the invariant field.
If \(q\) is the quotient by \(\tau\), the pullback and norm identities
on Jacobians give
\begin{equation}\label{d2:eq:klein-identity}
 q^*q_*+h_+^*h_{+*}+h_-^*h_{-*}
   =2\Dtwoid+\sum_{g\in\{1,\tau,\iota,\tau\iota\}}g
   =2\Dtwoid.
\end{equation}
The full group norm is zero because it factors through the Jacobian
of the rational total quotient.
The class \(P_1-P_{-1}\) is killed by \(q_*\).
If its images under both \(h_+\) and \(h_-\) were torsion,
\eqref{d2:eq:klein-identity} would make it torsion.
Choose a quotient \(h\) for which the image is nontorsion.
Since \(h\tau=[-1]h\), this image is \(2Q\), where \(Q=h(P_1)\).
Every point in \(E_0(k)\) is defined over a finite field and is
torsion. Hence \(Q\) is nonconstant.
The lifted root section gives
\(\Dtwowt Q\colon\DtwoSpec D\to\DtwocE^\circ\) with this reduction.

Let \(i\in W\) be the fourth root of unity reducing to \(2\).
The rotation
\([X:Y:Z]\mapsto[iX:iY:Z]\) fixes \(s\), sends the transformed
\(w\) to \(-iw\), sends \(z\) to \(-z\), and sends \(v_\pm\) to
\(-iv_\pm\).
Thus it induces \([-i]\) on the elliptic target.
It sends \(P_1\) to \(P_i\), so
\begin{equation}\label{d2:eq:root-equivariance}
 h(P_g)=[g^{-1}]Q
\end{equation}
for all fourth roots \(g\), including the lifted Teichm\"uller
sections.

We need an integral trace on total forms, including base
differentials.
For a tame double quotient of smooth relative curves, the trace is
the \'etale trace away from ramification.
At ramification, choose an anti-invariant fiber parameter \(t\);
\'etale-locally the quotient parameter is \(s=t^2\).
Invariant fiber one-forms are multiples of
\(t\,dt=ds/2\).
Invariant base forms and their wedges also descend.
Thus invariant total forms are exactly the forms on the quotient.
Summing \(1+\iota^*\) and descending defines a chain map
\[
 \DtwoTr_hh^*=2,\qquad h^*\DtwoTr_h=1+\iota^*.
\]
It is integral and continuous, and extends through localization and
completion, first relative to \(W\) and then relative to \(R\).
For an invertible module \(L\), the Cartier-divisor norm identity is
\(h^*\DtwoNm_hL=L\otimes\iota^*L\), including ramification
multiplicities.
Since \(h^*\) is injective on rational cohomology by
\(\DtwoTr_hh^*=2\), it follows that
\begin{equation}\label{d2:eq:trace-chern}
 \DtwoTr_h(c_1^\DtwodR(L))=c_1^\DtwodR(\DtwoNm_hL).
\end{equation}
This is an identity on ordinary de Rham quotients, because the trace
is an actual chain map.

Consider first
\[
 B_{\Dtwocyc}=
 C'\langle X,Y\rangle/(X^4+Y^4+TXY^3-1),
 \qquad
 \gamma_g^\Dtwocyc=c_1^\DtwodR(\DtwocO(P_g)).
\]
The norm of its root line under \(h\) is the graph line of
\([g^{-1}]\Dtwowt Q\).
Choose \(\eta,\omega\) with the polarization convention of
\eqref{d2:eq:graph-class}.
The unit-root character can be determined exactly.
On the elliptic curve,
\[
 [i](z,v)=(-z,iv),\qquad
 [i]^*\!\left(\frac{dz}{2v}\right)
       =i\,\frac{dz}{2v}.
\]
By the perfect pairing, the character on \(H^1(\DtwocO)\) is
\(i^{-1}\).
A primitive integral generator of the unit-root crystalline line
has nonzero image in \(H^1(\DtwocO)\) modulo five.
Otherwise its nonzero reduction would lie in the holomorphic
subspace, which crystalline Frobenius kills modulo five,
contradicting its unit Frobenius eigenvalue.
The two order-four eigenvalues have distinct reductions.
Consequently the character on the unit-root line is exactly
\(\omega_T^3\), rather than an unspecified odd character.

Let \(\lambda_g=\sigma_g(b)\) for any anti-invariant \(b\).
Then \(\lambda_{-g}=-\lambda_g\).
Pair opposite roots and apply \eqref{d2:eq:trace-chern}, followed by
the contraction \(L\) of Proposition~\ref{d2:prop:contraction}.
The graph-difference formula \eqref{d2:eq:graph-class} contributes
twice the pullback of \(\eta\).
Equation~\eqref{d2:eq:root-equivariance} gives
\([g^{-1}]^*\eta=\omega_T(g)^{-3}\eta\).
Hence the exact linear identity is
\begin{equation}\label{d2:eq:linear-weight}
 L\DtwoTr_h\left(\sum_{g\in G}\sigma_g(b)\gamma_g^\Dtwocyc\right)
   =
 \left(\sum_{g\in G}\sigma_g(b)\omega_T(g)^{-3}\right)
       [\Dtwowt Q^*\eta]
   =4e_3(b)[\Dtwowt Q^*\eta].
\end{equation}
The factor two from the graph difference is exactly the contribution
of the two opposite roots in the full sum.
The polarization convention fixes the common sign in this formula.

Proposition~\ref{d2:prop:graph}, using
Lemma~\ref{d2:lem:scalar}, shows that \([\Dtwowt Q^*\eta]\) is nonzero
in the ordinary Cohen one-form quotient.
Its finite-constant-extension assertion preserves this class over
\(R\).
Thus \eqref{d2:eq:linear-weight} is nonzero when \(e_3(b)\ne0\);
in particular it is nonzero under the two nonvanishing assumptions
of the proposition.
Both operators preserve exactness, so the weighted class on
\(B_{\Dtwocyc}\) is nonzero.
Formula~\eqref{d2:eq:linear-weight} remains valid for all anti-invariant \(b\), even if an odd component vanishes.

Finally we transport all four root lines to \(B\).
The cyclotomic identity gives \(5/\pi^4\equiv-1\bmod\pi\), and all
other coefficients below degree four in
\(\Phi_5(1+\pi Y)/\pi^4\) are divisible by \(\pi\).
Thus \(B\) and \(B_{\Dtwocyc}\) have the same special fiber
\[
 X^4+Y^4+TXY^3=1.
\]
For \(G_0=X^4+Y^4+TXY^3\), Euler's identity
\(X(G_0)_X+Y(G_0)_Y=4G_0\) proves smoothness on this fiber.
The two equations are monic in \(Y\), so the models are
\(\pi\)-torsion-free. Their completions are formally smooth over
\(C'\).

Lift the identity modulo \(\pi\) successively modulo \(\pi^n\).
This gives a continuous \(C'\)-linear map between the two completed
models.
Successive approximation proves surjectivity.
An element of its kernel lies in \(\pi\) times the source;
torsion-freeness of the target permits division of the relation by
\(\pi\). Iteration and separatedness prove injectivity.
Thus the map is an isomorphism.

On the support of \(J_g\), the \(Y\)-derivative is
\(\prod_{h\ne g}(a_g-a_h)\), a unit.
The equation expresses \(Y-a_g\) locally as a multiple of \(X\);
elsewhere \(J_g\) is the unit ideal.
Hence \(J_g\) is invertible.
Its reduction is \((X,Y-g)\), equal to the reduction of the cyclic
root ideal.
In a complete ring, \(\pi\) lies in the Jacobson radical.
An isomorphism between the reductions of two projective rank-one
modules lifts by projectivity; Nakayama at every maximal ideal
makes the lift an isomorphism.
Therefore the formal replacement identifies every root Picard class.

The replacement fixes \(C'\), so it fixes all \(\sigma_g(b)\).
It induces an isomorphism of continuous forms relative to \(R\),
with \(dT\) retained.
The weighted nonzero class, and the transported linear functional
\eqref{d2:eq:linear-weight}, therefore give the assertion for \(B\).
\end{proof}

\section{An actual arithmetic coefficient}\label{d2:sec:arithmetic}

The goal of this section is to construct one actual coefficient whose
two odd point-character components are nonzero.

\begin{proposition}\label{d2:prop:coefficient}
Put \(W_E=\DtwoZZ_5[\rho]\), \(E=\DtwoQQ_5(\rho)\), and
\(\pi=\rho-1\).
Define
\[
 r_E\colon K_3(W_E)\longrightarrow E
\]
by passage to \(\pi_3\DtwocK(W_E)\), the inverse of
\(\pi_3\DtwocK(W_E,(\pi))\to\pi_3\DtwocK(W_E)\), and the relative point
equivalence with \(\DtwoHC_2^\Dtwocts(W_E/W_E;\DtwoQQ_5)=E\).
There is an actual class
\[
 \beta=(1-\sigma_{-1})\beta_0\in K_3(W_{\Dtwoar})
\]
whose reduction in \(K_3(\DtwoFF_5)\) is zero, such that
\(\sigma_{-1}\beta=-\beta\) integrally and
\[
 b=r_E(\beta|_{W_E})
 \quad\text{satisfies}\quad e_1b\ne0,\qquad e_3b\ne0.
\]
The projectors \(e_r\) act only on the indicated
\(\DtwoQQ_5\)-vector-space targets.
\end{proposition}

\begin{proof}
For \(r=1,3\), set
\[
 \tau_r=\sum_{g=1}^4\omega_T(g)^{-r}\rho^g.
\]
Its \(\pi\)-adic valuation is \(r\).
To see this, expand \(\rho^g=(1+\pi)^g\).
For \(j<r\), the coefficient
\(\sum_g\omega_T(g)^{-r}\binom gj\) vanishes modulo five:
\(\binom gj\) has degree \(j\) in \(g\), and the power sums on
\(\DtwoFF_5^\times\) vanish unless the exponent is divisible by four.
For \(j=r\), the coefficient is \(4/r!\bmod5\), which is nonzero.
Since \(5\) is a unit times \(\pi^4\), terms with \(j<r\)
have valuation at least \(4+j>r\), while terms with \(j>r\)
have valuation at least \(j>r\).
The degree-\(r\) term is the unique one of minimal valuation.

Let \(\DtwoLi_2\) be the Coleman dilogarithm normalized by
\(\sum_{m\ge1}z^m/m^2\) on the open unit disc.
Define
\[
 S_j(z)=\left(z\frac d{dz}\right)^j\frac1{1-z},
 \qquad
 F_2(z)=\DtwoLi_2(z)-5^{-2}\DtwoLi_2(z^5).
\]
Partitioning the indices prime to five into \(a+5m\), and expanding
\((a+5m)^{-2}\), gives on that disc
\begin{equation}\label{d2:eq:depletion}
 F_2(z)=
 \sum_{j\ge0}(-1)^j(j+1)5^jS_j(z^5)
       \sum_{a=1}^4\frac{z^a}{a^{j+2}}.
\end{equation}
The right side converges uniformly on
\(|z|\le1,\ |1-z^5|=1\).
Each \(S_j\) has integral coefficients and unit denominator there,
so its \(j\)-th term has norm at most \(5^{-j}\).
This affinoid is connected: its reduction is the integral curve
\(\DtwoSpec\DtwoFF_5[z,(1-z)^{-1}]\), and its integral affinoid algebra
has no nontrivial idempotent.

By \cite[Theorem~2.7(3)]{BesserDeJeuAlgorithm}, the depleted polylogarithm
\(\DtwoLi_n(z)-p^{-n}\DtwoLi_n(z^p)\) is analytic outside the closed disc
\(|z-1|\le p^{-1/(p-1)}\), with a convergent expression in
\(1/(1-z)\).
For \(p=5\), this region contains the affinoid just used.
Both sides of \eqref{d2:eq:depletion} are analytic there.
Their equality on the open unit disc extends by the identity
principle.

At \(z=-\rho\), \(z^5=-1\).
The value \(\DtwoLi_2(-1)\) is fixed by \(G\), so its odd projections
vanish.
Permuting the summation indices gives the exact identity
\[
 e_r(\rho^a)=\frac{\omega_T(a)^r}{4}\tau_r.
\]
Applying it to \eqref{d2:eq:depletion} yields
\begin{equation}\label{d2:eq:odd-dilogarithms}
 \frac{4e_r\DtwoLi_2(-\rho)}{\tau_r}
 =
 \sum_{j\ge0}(-1)^j(j+1)5^jS_j(-1)
    \sum_{a=1}^4\frac{(-1)^a\omega_T(a)^r}{a^{j+2}}.
\end{equation}
All \(S_j(-1)\) are \(5\)-integral.
Modulo five, only the \(j=0\) term remains:
\[
 \frac12\sum_{a=1}^4(-1)^aa^{r-2}.
\]
For \(r=1\) it is
\(\frac12(-1+3-2+4)=2\) in \(\DtwoFF_5\).
For \(r=3\) it is
\(\frac12(-1+2-3+4)=1\).
Thus both odd projections of \(\DtwoLi_2(-\rho)\) are nonzero.
Division by \(\tau_r\) is legitimate by its valuation, and the
projection identity is exact, so no precision loss is involved.

The special-unit symbol construction
\cite[Theorems~1.6 and~1.10, Remark~1.7]{BesserDeJeu}
maps the degree-one, weight-two symbol complex of a localized
number ring to rational K-theory.
Its syntomic regulator is
\([x]_n\mapsto\pm(n-1)!L_{\mathrm{mod},n}(x)\);
the sign is the single convention in the cited remark.
In weight two there is no conjectural vanishing assumption.

Take \(x=-\rho\) in \(W_{\Dtwoar}\).
Both \(x\) and \(1-x=1+\rho\) are units, since the latter reduces
to \(2\) modulo \(\pi\).
The symbol \([-\rho]_2\) is a rational degree-one cocycle:
its differential is \((1+\rho)\otimes(-\rho)\), or the exterior
version, and the root-of-unity factor is zero after rationalization.
It gives \(\beta_{\mathrm{rat}}\in K_3(W_{\Dtwoar})\otimes\DtwoQQ\).
Its syntomic regulator is \(\pm\DtwoLi_2(-\rho)\), because the
modification terms contain \(\log(-\rho)=0\).
Choose a positive integer \(N\) and an actual \(\beta_0\)
representing \(N\beta_{\mathrm{rat}}\), and set
\(\beta=(1-\sigma_{-1})\beta_0\).
This is an operation on actual groups.
The automorphism \(\sigma_{-1}\) induces the identity on the residue
field, so \(\beta\) reduces to zero.
Its integral antisymmetry follows by squaring \(\sigma_{-1}\).
Equivariance of the regulator and
\eqref{d2:eq:odd-dilogarithms} show that its two odd syntomic
components are nonzero.

For the possibly ramified field \(E\), the syntomic-to-\'etale
comparison in \cite[Section~2]{HuberKings} has no ramification
restriction.
The degree-one syntomic coordinate for weight \(n>1\) is \(E\)
\cite[Example~2.2.4]{HuberKings}.
It agrees with the definition used above
\cite[Proposition~2.2.7]{HuberKings}.
Its Chern-compatible map to \(H^1(E,\DtwoQQ_5(n))\) is the
Bloch--Kato exponential
\cite[Propositions~2.2.9 and~2.3.4]{HuberKings}, which is an
isomorphism for \(n>1\), as recalled in
\cite[Section~1.3]{HuberKings}.
These maps are natural under field automorphisms.
Thus the actual \'etale Chern class
\(c_{2,1}(\beta)\in H^1(E,\DtwoQQ_5(2))\) has both odd components
nonzero.

There is a natural \(G\)-equivariant \(\DtwoQQ_5\)-linear map
\begin{equation}\label{d2:eq:completed-chern}
 \pi_3\DtwocK(W_E)\longrightarrow H^1(E,\DtwoQQ_5(2))
\end{equation}
through which the Chern class of every actual element factors.
To construct this map, first work before inverting five.
Map from \(\pi_3(K(W_E)^\wedge_5)\) to \(\pi_3(K(E)^\wedge_5)\), and then to
\(\varprojlim_nK_3(E;\DtwoZZ/5^n)\).
Apply the finite-coefficient Chern classes
\[
 K_3(E;\DtwoZZ/5^n)\longrightarrow
 H^1(E,\mu_{5^n}^{\otimes2}).
\]
Naturality, compatibility with reduction of integral classes,
additivity for odd coefficients, and compatibility with coefficient
transitions are respectively
\cite[Proposition~2.1.1, Lemma~2.3, Propositions~2.4 and~2.8]{WeibelChern}.
The resulting compatible additive maps recover the actual Chern
class. Each finite target is killed by \(5^n\), so the inverse-limit
map is \(\DtwoZZ_5\)-linear.
The groups \(H^0(E,\mu_{5^n}^{\otimes2})\) are finite, hence
Mittag--Leffler. The continuous-cohomology inverse-limit sequence
therefore identifies the target limit with \(H^1(E,\DtwoZZ_5(2))\).
Inverting five gives \eqref{d2:eq:completed-chern}.
No derived inverse-limit vanishing on the K-theory side is needed.

The finite-field calculation \cite{QuillenFinite}\cite[IV, Corollary~1.13]{Weibel} gives
\[
 K_{2j}(\DtwoFF_5)=0,\qquad
 K_{2j-1}(\DtwoFF_5)=\DtwoZZ/(5^j-1)\qquad(j\ge1).
\]
Multiplication by five is invertible on these positive groups.
The finite-coefficient exact sequences show that all positive
homotopy groups of \(K(\DtwoFF_5)^\wedge_5\) vanish.
In particular the groups in degrees three and four of
\(\DtwocK(\DtwoFF_5)\) vanish.
The relative fiber sequence identifies
\(\pi_3\DtwocK(W_E,(\pi))\) with \(\pi_3\DtwocK(W_E)\).

The point case of
\cite[Theorem~A, Proposition~4.10, Corollary~3.8]{AMMN}
identifies this group with
\(\DtwoHC_2^\Dtwocts(W_E/W_E;\DtwoQQ_5)=E\).
Here \(W_E\) is a complete mixed-characteristic discrete valuation
ring with perfect finite residue field, its formal spectrum is
smooth over itself, and changing \((5)\) to \((\pi)\) changes the
special fiber only by a nilpotent thickening.
The comparison is natural for automorphisms of the coefficient
point, hence \(G\)-equivariant.

Transporting \eqref{d2:eq:completed-chern} through this point
identification gives a \(G\)-equivariant \(\DtwoQQ_5\)-linear map from
\(E\) to \(H^1(E,\DtwoQQ_5(2))\), carrying \(b\) to
\(c_{2,1}(\beta)\).
If \(e_rb=0\), the corresponding Chern component would vanish.
This contradicts the syntomic calculation for \(r=1,3\).
Equivariance also gives \(\sigma_{-1}b=-b\).
No equality between a chosen syntomic normalization and the point
coordinate was required.
\end{proof}

\section{Assembly over the local ring and its completion}\label{d2:sec:assembly}

In this section, we prove Theorem~\ref{d2:thm:completed} and show that its
actual coefficient also gives the geometric detector.

\begin{proof}[Proof of Theorem~\ref{d2:thm:completed}]
Choose the actual \(\beta\) of Proposition~\ref{d2:prop:coefficient},
and set \(c=c(\beta)\).
Section~\ref{d2:sec:preliminaries} proves regularity of \(A,\Dtwowh A\)
and both integral generic vanishings.
We prove that the image of this same \(c\) in \(K_3(\Dtwowh A)\) has
infinite order.

Put \(k=\overline{\DtwoFF}_5\), \(W=W(k)\), and \(R=W[\rho]\).
Choose \(C\), arising from a smooth affine curve \(D\), as in
Proposition~\ref{d2:prop:weighted}.
Write \(L\) for the function field of \(D/5D\).
It is a finite separable extension of \(k(T)\).
The coefficient ring \(C'=C\otimes_WR\) is a complete discrete
valuation ring with uniformizer \(\pi\) and residue field \(L\).
Use the test algebra
\begin{equation}\label{d2:eq:test-algebra}
 B=C'\langle X,Y\rangle/
 \left(X^4+\frac{\Phi_5(1+\pi Y)}{\pi^4}+TXY^3\right).
\end{equation}
The polynomial is integral over \(R\): every coefficient below
degree four before division has \(\pi\)-valuation at least four.
It is monic in \(Y\), so the polynomial quotient and its completion
are \(\pi\)-torsion-free.

Its special fiber is
\begin{equation}\label{d2:eq:test-special}
 B/\pi B=L[X,Y]/(X^4+Y^4+TXY^3-1).
\end{equation}
For \(G_0=X^4+Y^4+TXY^3\), Euler's identity gives
\(X(G_0)_X+Y(G_0)_Y=4G_0=4\) on this fiber.
It is therefore smooth over \(L\).
Its projective closure is the smooth geometrically integral quartic
of Proposition~\ref{d2:prop:weighted}, so the fiber is a domain.

The ring \(B\) is noetherian by its monic presentation and adic
completion. Since \(\pi\) lies in its Jacobson radical, all maximal
ideals contain \(\pi\).
At each one, the quotient by \(\pi\) is regular and \(\pi\) is a
nonzero divisor. Lifting a regular system of parameters from the
quotient and adjoining \(\pi\) proves regularity of that local ring.
Hence \(B\) is regular.
Its dimension is two: the restricted two-variable ring over \(C'\)
has dimension three, and the equation is a nonzero monic equation.
Equivalently, the parameter argument gives dimension at most two
at every maximal ideal and equality at closed points of the special
curve.
Separatedness, torsion-freeness, and the domain special fiber show
that \(B\) is a domain: the first nonzero \(\pi\)-adic terms of a
product of two nonzero elements have nonzero product.

We construct the essentially smooth rank-two model required in
Proposition~\ref{d2:prop:cyclic}.
Let \(U_D=D_{5D}\), whose completion is \(C\), and put
\[
 H=
 \bigl((U_D\otimes_WR)[X,Y]\bigr)/
 \left(X^4+\frac{\Phi_5(1+\pi Y)}{\pi^4}+TXY^3\right).
\]
Its completion is \(B\).
The two relative partial derivatives generate the unit ideal
modulo \(\pi\) by the Euler calculation.
Choose an element \(j\) of their generated ideal with
\(j\equiv1\bmod\pi\).
Then \(H[1/j]\) is smooth of relative dimension one over
\(U_D\otimes_WR\).
Inverting \(j\) does not change the completion.

The ring \(U_D\otimes_WR\) is a localization of the smooth
\(R\)-algebra \(D\otimes_WR\), of relative dimension one.
Transitivity of smoothness makes \(H[1/j]\) a localization of a
smooth \(R\)-algebra with relative cotangent rank two.
The finite coefficients and the choice of \(j\) spread over a
further localization of \(D\).
This gives the required \(B_0\).
The \(5\)-adic and \(\pi\)-adic completions agree because \(5\)
is a unit times \(\pi^4\).
In particular, relative forms over \(R\) retain the differential
\(dT\).

Write \(b=r_E(\beta|_{W_E})\).
Proposition~\ref{d2:prop:weighted} gives invertible ideals
\[
 J_g=\left(X,Y-\frac{\rho^g-1}{\pi}\right)
\]
and the nonzero class
\begin{equation}\label{d2:eq:completed-form}
 \Xi=\sum_{g\in G}\sigma_g(b)\bigl(-c_1^\DtwodR(J_g)\bigr)
 \ne0
 \quad\text{in}\quad
 \Omega^2_{B/R}[1/5]\big/d\Omega^1_{B/R}[1/5].
\end{equation}

The reduction of \(\beta|_{W_E}\) in \(K_3(\DtwoFF_5)\) is zero.
It lifts to an actual
\(\beta_{\mathrm{rel}}\in K_3(W_E,(\pi))\).
This lift is unique, since \(K_4(\DtwoFF_5)=0\) by
\cite{QuillenFinite}\cite[IV, Corollary~1.13]{Weibel}.
Naturality makes the lifts of all conjugates compatible.
Pull them through \(W_E\to R\).
The relative point coordinates are \(\sigma_g(b)\) in \(R[1/5]\):
the point comparison is natural for this coefficient extension,
and the map on degree-two point cyclic homology is inclusion of
the coefficient fields.

Form the actual relative class
\[
 \delta_{\mathrm{rel}}
   =\sum_{g\in G}
      (\sigma_g\beta_{\mathrm{rel}})|_B\cdot(1-[J_g])
       \in K_3(B,\pi B),
\]
and let \(\delta\) be its image in actual \(K_3(B)\).
Proposition~\ref{d2:prop:cyclic}, applied to the rank-two model above,
identifies its top-form component with \(\lambda_R\Xi\), where
\(\lambda_R\ne0\).
Thus the image of \(\delta_{\mathrm{rel}}\) in
\(\pi_3\DtwocK(B,\pi B)\) is nonzero.

We establish the special-fiber vanishing before using either fiber
sequence.
The ring in \eqref{d2:eq:test-special} is the filtered colimit of
\[
 (D/5D)[1/s][X,Y]/(X^4+Y^4+TXY^3-1),
 \qquad 0\ne s\in D/5D.
\]
These are smooth affine surfaces over the perfect field \(k\).
The base is a smooth affine curve, and the Euler identity makes
the displayed curve over it smooth.
For a smooth \(d\)-dimensional variety over a perfect field of
characteristic \(p\), the finite-coefficient theorem
\cite[Theorem~8.4]{GeisserLevineCharP} gives
\(K_j(-;\DtwoZZ/p^n)=0\) for \(j>d\).
Its groups are the homotopy groups of the cofiber spectrum, with
their torsion contribution, rather than mere tensor products of
actual K-groups.
The theorem follows there from the motivic-to-K spectral sequence
and the vanishing of logarithmic de Rham--Witt sheaves above the
dimension.
Applying its smooth-surface form and then filtered continuity gives
\[
 K_j(B/\pi B;\DtwoZZ/5^n)=0
 \qquad(j>2,\ n\ge1).
\]
For \(j\ge3\), both the degree-\(j\) and degree-\((j+1)\) groups
of the finite-coefficient tower vanish.
The Milnor exact sequence therefore gives
\(\pi_jK(B/\pi B)^\wedge_5=0\), and inversion of five preserves
this vanishing. In particular,
\begin{equation}\label{d2:eq:special-vanishing}
 \pi_3\DtwocK(B/\pi B)=\pi_4\DtwocK(B/\pi B)=0.
\end{equation}

First use the reduction fiber sequence
\begin{equation}\label{d2:eq:reduction-fiber}
 \DtwocK(B,\pi B)\longrightarrow\DtwocK(B)
                  \longrightarrow\DtwocK(B/\pi B).
\end{equation}
Equation~\eqref{d2:eq:special-vanishing} makes the first map an
isomorphism on degree three.
The image of \(\delta\) in \(\pi_3\DtwocK(B)\) is consequently nonzero.

Next use regularity of \(B\), regularity of \(B/\pi B\), and
the nonzero divisor \(\pi\).
Localization and d\'evissage
\cite[V.4.1, V.6.1]{Weibel} give the different fiber sequence
\begin{equation}\label{d2:eq:localization-fiber}
 \DtwocK(B/\pi B)\longrightarrow\DtwocK(B)
                         \longrightarrow\DtwocK(B[1/\pi]).
\end{equation}
The left map is support pushforward, not reduction.
Its source has zero degree-three group, so the second map is
injective on degree three.
Thus \(\delta|_{B[1/\pi]}\) has nonzero image in
\(\pi_3\DtwocK(B[1/\pi])\).
That target is a \(\DtwoQQ_5\)-vector space, so the actual class
\(\delta|_{B[1/\pi]}\) has infinite order.
Only the fact that torsion maps to zero in a characteristic-zero
vector space is used here.

There is a ring map \(A\to B\) defined by
\[
 x\longmapsto\pi X,\qquad y\longmapsto\pi Y,
 \qquad T\longmapsto T.
\]
The defining equation maps to \(\pi^4\) times the equation of \(B\).
Every element inverted in forming \(A\) reduces to a nonzero
polynomial in \(T\) modulo \((5,x,y)\).
Its image modulo \(\pi\) is a nonzero element of \(L\), hence a
unit. Since \(\pi\) lies in the Jacobson radical, its image is
a unit in \(B\).
This verifies the map on the localization.
It takes \((x,y)\) into \(\pi B\), so completeness extends it
continuously to \(\Dtwowh A\to B\).

We match the actual supported class only after a flat splitting.
Set
\[
 S_{\mathrm{spl}}=
 (A\otimes_{\DtwoZZ_{(5)}}W_{\Dtwoar})[1/5].
\]
This is flat over \(A\), and regular because it is finite \'etale
over \(A[1/5]\).
The polynomial \(P(y)\) splits there with roots \(\rho^g-1\),
whose differences are units.
The divisor \(x=0\) is the disjoint union of the Cartier divisors
with ideals
\[
 I_g=(x,y-(\rho^g-1)).
\]
Near a component, the other factors of \(P\) are units, and the
equation expresses \(y-(\rho^g-1)\) as a multiple of \(x\).
Away from that component the ideal is the unit ideal.
The element \(x\) remains a nonzero divisor by flatness, and the
component rings are localizations of \(E_{\Dtwoar}[T]\), hence regular.

On the \(g\)-th component, the original coefficient map sends
\(\rho=1+y\) to \(\rho^g\).
Flat base change and the projection formula give the equality of
actual classes
\begin{equation}\label{d2:eq:flat-splitting}
 c|_{S_{\mathrm{spl}}}
   =\sum_{g\in G}(\sigma_g\beta)|_{S_{\mathrm{spl}}}
                         \cdot(1-[I_g]).
\end{equation}
No scalar projector acts on an actual K-group in this identity.

The map \(A\to B\), together with the coefficient embedding,
gives \(S_{\mathrm{spl}}\to B[1/\pi]\).
There
\[
 x=\pi X,\qquad
 y-(\rho^g-1)=\pi(Y-a_g).
\]
The pulled-back invertible module \(I_g\) identifies with
\(J_g[1/\pi]\).
Locally on its support it is generated by \(x\), whose image
\(\pi X\) is a nonzero divisor in the domain \(B[1/\pi]\);
off the support it is generated by a unit.
Thus this identification is an isomorphism of invertible modules,
not only an equality of generated ideals.

Pulling back \eqref{d2:eq:flat-splitting} now gives
\[
 c|_{B[1/\pi]}
 =\sum_g(\sigma_g\beta)|_{B[1/\pi]}
                          (1-[J_g[1/\pi]])
 =\delta|_{B[1/\pi]}.
\]
The right side has infinite order.
The map factors through \(\Dtwowh A\), so it is also the image of
\(\Dtwowh c\). Hence \(\Dtwowh c\), and therefore \(c\), has infinite order.

An element of an abelian group maps to zero after tensoring with
\(\DtwoQQ\) precisely when it is torsion.
The two integral generic vanishings established at the start thus
give the two asserted nonzero rational kernels.

We finish by identifying this coefficient with one allowed in the
geometric proof. The inclusions
\(\DtwoZZ[\rho]\subset\DtwocO_{E_{\Dtwoar}}\subset W_{\Dtwoar}\) show that
\(W_{\Dtwoar}\) is the localization of \(\DtwocO_{E_{\Dtwoar}}\) obtained by
inverting integers prime to five.
Localization, filtered continuity, and
\(K_2\) of finite fields being zero give a surjection
\[
 K_3(\DtwocO_{E_{\Dtwoar}})\longrightarrow K_3(W_{\Dtwoar});
\]
see \cite[V.6.1]{Weibel},
\cite[IV, Section~6.4]{Weibel}, and
\cite{QuillenFinite}\cite[IV, Corollary~1.13]{Weibel}.
Lift the actual \(\beta_0\) in
Proposition~\ref{d2:prop:coefficient} to
\(\gamma_0\in K_3(\DtwocO_{E_{\Dtwoar}})\).
Then \(\gamma=(1-\sigma_{-1})\gamma_0\) maps exactly to \(\beta\).

The nonzero point coordinate makes \(\beta\) nontorsion, so
\(\gamma\) is nontorsion.
The regulator injectivity of
\cite[Section~6]{BunkeTamme}, with the rank calculation of
\cite[Proposition~12.2]{Borel}, gives \(\gamma\) a nonzero
complex regulator.
All calculations in Section~\ref{d2:sec:analytic} therefore apply to
this same \(\beta\).
Thus both detectors apply to the actual class \(c\) just constructed,
and its image under completion is exactly \(\Dtwowh c\).
\end{proof}
